% Generado por la construcción de los notebooks (no editar a mano).
% Índice de macros y tablas por etapa. % Generado por la construcción de los notebooks (no editar a mano).
% Cifras citadas en el informe; el índice comentado está en indice_macros.tex.

% --- 01_datos_exploracion ---
\newcommand{\DatBloquesCV}{4}
\newcommand{\DatClientes}{220}
\newcommand{\DatClientesConMuestras}{219}
\newcommand{\DatColumnasContrato}{65}
\newcommand{\DatColumnasObligatorias}{15}
\newcommand{\DatColumnasOpcionales}{50}
\newcommand{\DatControles}{26236}
\newcommand{\DatCortePrueba}{2026-01-01}
\newcommand{\DatCorteValidacion}{2025-07-01}
\newcommand{\DatCostoFN}{5}
\newcommand{\DatCostoFP}{1}
\newcommand{\DatDesde}{2023-01-02}
\newcommand{\DatDiasCalendario}{1398}
\newcommand{\DatDuplicadas}{148}
\newcommand{\DatFaltanOpcionales}{0}
\newcommand{\DatFechaFin}{2026-10-30}
\newcommand{\DatFechaInicio}{2023-01-02}
\newcommand{\DatHasta}{2026-10-30}
\newcommand{\DatHuella}{f82dd7cc9bcb}
\newcommand{\DatMomentos}{9}
\newcommand{\DatMuestrasCrudas}{38994}
\newcommand{\DatOrigen}{simulacion}
\newcommand{\DatRecuentosNoNumericos}{28716}
\newcommand{\DatRecuentosNoNumericosPct}{75.9}
\newcommand{\DatSemilla}{42}
\newcommand{\DatTablasContrato}{4}
\newcommand{\DatValoresImposiblesCrudos}{366}
\newcommand{\DatValoresImposiblesEnDuplicadas}{1}
\newcommand{\DatVisitas}{1812}
\newcommand{\ExpCVSemanal}{0.30}
\newcommand{\ExpCategoriaMasNiveles}{codigo\_producto}
\newcommand{\ExpClientesDesarrollo}{203}
\newcommand{\ExpClientesMitad}{28}
\newcommand{\ExpClientesMitadPct}{13.8}
\newcommand{\ExpColumnasConFaltantes}{25}
\newcommand{\ExpColumnasIntegradas}{69}
\newcommand{\ExpColumnasSoloEstructural}{17}
\newcommand{\ExpConformesPorNoConforme}{11}
\newcommand{\ExpDecimoClientesPct}{41.1}
\newcommand{\ExpDiferenciaEstandarMax}{0.34}
\newcommand{\ExpDuplicadosEliminados}{148}
\newcommand{\ExpEnsayoMasTardio}{mohos y levaduras}
\newcommand{\ExpEnsayoMasTardioPct}{16.6}
\newcommand{\ExpEnsayoMenosTardio}{Staphylococcus aureus}
\newcommand{\ExpEnsayoMenosTardioPct}{8.8}
\newcommand{\ExpEntregaMediana}{6}
\newcommand{\ExpEntregaPNoventa}{12}
\newcommand{\ExpEntregaTardiaPct}{12.6}
\newcommand{\ExpFilasEntrenamiento}{23151}
\newcommand{\ExpFilasPrueba}{8947}
\newcommand{\ExpFilasValidacion}{6748}
\newcommand{\ExpHorasCuartoAltoPct}{12.6}
\newcommand{\ExpHorasCuartoBajoPct}{7.0}
\newcommand{\ExpIndiceMesAlto}{1.26}
\newcommand{\ExpIndiceMesBajo}{0.83}
\newcommand{\ExpMedianaMuestrasCliente}{73}
\newcommand{\ExpMesAlto}{nov}
\newcommand{\ExpMesBajo}{ene}
\newcommand{\ExpMuestrasCrudas}{38994}
\newcommand{\ExpMuestrasDesarrollo}{29899}
\newcommand{\ExpMuestrasLimpias}{38846}
\newcommand{\ExpMuestrasSemana}{192}
\newcommand{\ExpMuestrasSemanaMax}{366}
\newcommand{\ExpMuestrasSemanaMin}{61}
\newcommand{\ExpNivelesCategoria}{1169}
\newcommand{\ExpParCorrelacionado}{dias\_prometidos y limite\_ufc}
\newcommand{\ExpPeorFaltanteColumna}{actividad\_agua}
\newcommand{\ExpPeorFaltantePct}{9.8}
\newcommand{\ExpPrimerAnio}{2023}
\newcommand{\ExpProblemaMayor}{recuento censurado (``<N'' o ``>N'')}
\newcommand{\ExpProblemaMayorFilas}{18816}
\newcommand{\ExpReglasConHallazgos}{18}
\newcommand{\ExpRhoObjetivo}{0.07}
\newcommand{\ExpRhoPar}{$-$0.28}
\newcommand{\ExpSectorMayor}{cosmético}
\newcommand{\ExpSectorMayorPct}{36.3}
\newcommand{\ExpSectorMenor}{central de mezclas}
\newcommand{\ExpSectorMenorPct}{8.6}
\newcommand{\ExpSemanas}{157}
\newcommand{\ExpTablasLimpias}{5}
\newcommand{\ExpTasaMensualMaxPct}{12.7}
\newcommand{\ExpTasaMensualMinPct}{5.5}
\newcommand{\ExpTasaNoConformePct}{8.6}
\newcommand{\ExpTasaPrimerAnioPct}{7.9}
\newcommand{\ExpTasaRechazoPct}{3.0}
\newcommand{\ExpTasaUltimoAnioPct}{8.7}
\newcommand{\ExpTemperaturaCuartoAltoPct}{12.8}
\newcommand{\ExpTemperaturaCuartoBajoPct}{6.9}
\newcommand{\ExpUltimoAnio}{2025}
\newcommand{\ExpVariableCorrObjetivo}{actividad\_agua}
\newcommand{\ExpVariableMasSepara}{horas\_transporte}

% --- 02_base_real ---
\newcommand{\RealArbolGiniAntes}{0.0690}
\newcommand{\RealArbolGiniDespues}{0.0666}
\newcommand{\RealArbolGiniReduccion}{0.0023}
\newcommand{\RealArbolPrimerCorte}{tasa\_cliente > 0.0735105}
\newcommand{\RealCalcIntercepto}{$-$3.677}
\newcommand{\RealCalcLogit}{2.047}
\newcommand{\RealCalcMayor}{tasa\_cliente}
\newcommand{\RealCalcMayorAporte}{5.442}
\newcommand{\RealCalcProba}{0.886}
\newcommand{\RealCalcTerminos}{9}
\newcommand{\RealCampanaMuestrasMediana}{14}
\newcommand{\RealCampanaUnDiaPct}{69}
\newcommand{\RealCampanas}{385}
\newcommand{\RealClientes}{82}
\newcommand{\RealClientesDispersion}{28}
\newcommand{\RealClientesPlanta}{8}
\newcommand{\RealClientesProducto}{22}
\newcommand{\RealCompuestoPlanta}{59.98}
\newcommand{\RealCompuestoProducto}{6.00}
\newcommand{\RealCongeladoHuella}{10900224adf7}
\newcommand{\RealCongeladoModelo}{Regresión logística (C = 0.01)}
\newcommand{\RealCostoFN}{5}
\newcommand{\RealCostoFP}{1}
\newcommand{\RealDesde}{2024-06-04}
\newcommand{\RealDispLotesPlanta}{1.20}
\newcommand{\RealDispLotesProducto}{0.73}
\newcommand{\RealDispLotesTodos}{0.90}
\newcommand{\RealDispMuestrasMax}{407}
\newcommand{\RealDispMuestrasPlanta}{71.16}
\newcommand{\RealDispMuestrasProducto}{4.25}
\newcommand{\RealDispMuestrasTodos}{7.23}
\newcommand{\RealFilas}{120114}
\newcommand{\RealHasta}{2025-12-19}
\newcommand{\RealHojaFilas}{51}
\newcommand{\RealHojaFraccion}{0.314}
\newcommand{\RealHojaVeces}{8.8}
\newcommand{\RealLoteDeltaGeoPlanta}{448}
\newcommand{\RealLoteDeltaGeoProducto}{5.4}
\newcommand{\RealLoteMaxPlanta}{675}
\newcommand{\RealLoteMaxProducto}{60}
\newcommand{\RealLoteMediaPlanta}{18.7}
\newcommand{\RealLoteMediaProducto}{4.1}
\newcommand{\RealLoteMedianaPlanta}{6}
\newcommand{\RealLoteMedianaProducto}{3}
\newcommand{\RealLoteMejorPlanta}{1 + binomial negativa}
\newcommand{\RealLoteMejorProducto}{1 + binomial negativa}
\newcommand{\RealLoteRPlanta}{0.39}
\newcommand{\RealLoteRProducto}{0.88}
\newcommand{\RealLotes}{2543}
\newcommand{\RealLotesPlanta}{762}
\newcommand{\RealLotesProducto}{1781}
\newcommand{\RealMinimoLotes}{20}
\newcommand{\RealMuestras}{21549}
\newcommand{\RealNivel}{0.05}
\newcommand{\RealNoConformePct}{3.7}
\newcommand{\RealOddsBaja}{previas\_lote}
\newcommand{\RealOddsBajaValor}{0.70}
\newcommand{\RealOddsSube}{tasa\_cliente}
\newcommand{\RealOddsSubeValor}{1.76}
\newcommand{\RealPlantaPct}{66.2}
\newcommand{\RealPruebaAP}{0.072}
\newcommand{\RealPruebaAPConocidos}{0.073}
\newcommand{\RealPruebaAPInf}{0.059}
\newcommand{\RealPruebaAPNuevos}{0.097}
\newcommand{\RealPruebaAPRegla}{0.064}
\newcommand{\RealPruebaAPReglaInf}{0.053}
\newcommand{\RealPruebaAPReglaSup}{0.081}
\newcommand{\RealPruebaAPSup}{0.091}
\newcommand{\RealPruebaAUC}{0.769}
\newcommand{\RealPruebaAUCInf}{0.743}
\newcommand{\RealPruebaAUCSup}{0.794}
\newcommand{\RealPruebaClientes}{72}
\newcommand{\RealPruebaClientesNuevos}{17}
\newcommand{\RealPruebaCosto}{1718}
\newcommand{\RealPruebaCostoNada}{1710}
\newcommand{\RealPruebaFN}{339}
\newcommand{\RealPruebaFP}{23}
\newcommand{\RealPruebaFUno}{0.016}
\newcommand{\RealPruebaMuestras}{16325}
\newcommand{\RealPruebaNuevasPct}{29}
\newcommand{\RealPruebaPositivos}{342}
\newcommand{\RealPruebaPrecision}{0.115}
\newcommand{\RealPruebaPrevalencia}{0.021}
\newcommand{\RealPruebaPrevalenciaPct}{2.1}
\newcommand{\RealPruebaRecall}{0.009}
\newcommand{\RealPruebaSinHistoriaPct}{2.2}
\newcommand{\RealPruebaTN}{15960}
\newcommand{\RealPruebaTP}{3}
\newcommand{\RealPruebaVeces}{3.4}
\newcommand{\RealRiesgoAPArbol}{0.073}
\newcommand{\RealRiesgoAPArbolProfundoEnt}{0.212}
\newcommand{\RealRiesgoAPArbolProfundoVal}{0.073}
\newcommand{\RealRiesgoAPBase}{0.041}
\newcommand{\RealRiesgoAPElegido}{0.109}
\newcommand{\RealRiesgoAPHistoria}{0.105}
\newcommand{\RealRiesgoAPLogistica}{0.109}
\newcommand{\RealRiesgoAPLogisticaEnt}{0.114}
\newcommand{\RealRiesgoAUCArbol}{0.683}
\newcommand{\RealRiesgoAUCElegido}{0.689}
\newcommand{\RealRiesgoAUCHistoria}{0.673}
\newcommand{\RealRiesgoAUCLogistica}{0.689}
\newcommand{\RealRiesgoArbol}{profundidad 5}
\newcommand{\RealRiesgoElegido}{Regresión logística (C = 0.01)}
\newcommand{\RealRiesgoEntrenamiento}{15628}
\newcommand{\RealRiesgoEntrenamientoPct}{73}
\newcommand{\RealRiesgoHojaEsperados}{1.8}
\newcommand{\RealRiesgoLogistica}{C = 0.01}
\newcommand{\RealRiesgoMinimoHoja}{50}
\newcommand{\RealRiesgoPrevalenciaEntPct}{3.6}
\newcommand{\RealRiesgoPrevalenciaPct}{4.1}
\newcommand{\RealRiesgoSinDictamen}{98}
\newcommand{\RealRiesgoValidacion}{5823}
\newcommand{\RealRiesgoValidacionPct}{27}
\newcommand{\RealRiesgoVariables}{9}
\newcommand{\RealSegClientes}{46}
\newcommand{\RealSegK}{4}
\newcommand{\RealSegMayor}{Más muestras por semana y más dispersión semanal}
\newcommand{\RealSegMayorMuestrasPct}{56.3}
\newcommand{\RealSegMinimoMuestras}{20}
\newcommand{\RealSegNombreCuatro}{Más muestras por semana y más dispersión semanal}
\newcommand{\RealSegNombreDos}{Más muestreo en planta y más tamaño de lote}
\newcommand{\RealSegNombreTres}{Menos muestreo en planta y menos lotes por semana}
\newcommand{\RealSegNombreUno}{Más semanas con envío y más lotes por semana}
\newcommand{\RealSegSilueta}{0.374}
\newcommand{\RealSimIndiceAlto}{3.7}
\newcommand{\RealSimIndiceBajo}{2.2}
\newcommand{\RealSobreLotesPlantaPct}{25}
\newcommand{\RealSobreLotesProductoPct}{5}
\newcommand{\RealSobreMuestrasPlantaPct}{100}
\newcommand{\RealSobreMuestrasProductoPct}{95}
\newcommand{\RealUmbral}{0.61}
\newcommand{\RealUmbralAlertas}{41}
\newcommand{\RealUmbralAlertasSemana}{1.6}
\newcommand{\RealUmbralCosto}{1152}
\newcommand{\RealUmbralCostoNada}{1195}
\newcommand{\RealUmbralCostoTodo}{5584}
\newcommand{\RealUmbralFN}{225}
\newcommand{\RealUmbralFP}{27}
\newcommand{\RealUmbralFUno}{0.100}
\newcommand{\RealUmbralPrecision}{0.341}
\newcommand{\RealUmbralRecall}{0.059}
\newcommand{\RealUmbralTN}{5557}
\newcommand{\RealUmbralTP}{14}
\newcommand{\RealUmbralTeorico}{0.167}
\newcommand{\RealValidacionVeces}{2.7}

% --- 03_cribado_regresion ---
\newcommand{\CriAgrCandidatos}{3}
\newcommand{\CriAgrComplementario}{segmentos\_clientes}
\newcommand{\CriAnalisisConTabla}{13}
\newcommand{\CriAuditoriasSuperadas}{13}
\newcommand{\CriBayes}{3}
\newcommand{\CriCandidatos}{16}
\newcommand{\CriCandidatosAgr}{3}
\newcommand{\CriCandidatosBayes}{3}
\newcommand{\CriCandidatosClf}{5}
\newcommand{\CriCandidatosReg}{5}
\newcommand{\CriClfMasHabil}{falso\_positivo\_lectura}
\newcommand{\CriClfMasHabilHabilidad}{0.689}
\newcommand{\CriClfMasHabilProporcion}{95}
\newcommand{\CriCribados}{13}
\newcommand{\CriDescartadosReglaSimple}{3}
\newcommand{\CriElegibles}{10}
\newcommand{\CriElegidos}{3}
\newcommand{\CriHopkinsMax}{0.86}
\newcommand{\CriHopkinsMin}{0.72}
\newcommand{\CriMasInestable}{cliente\_inactivo\_trimestre}
\newcommand{\CriMasInestableMax}{0.708}
\newcommand{\CriMasInestableMin}{0.442}
\newcommand{\CriMejorSilueta}{0.433}
\newcommand{\CriMejorSiluetaAnalisis}{perfiles\_muestra}
\newcommand{\CriMejorSiluetaK}{2}
\newcommand{\CriMinimoFilas}{300}
\newcommand{\CriMinimoFilasAgrupamiento}{50}
\newcommand{\CriMinimoPositivos}{100}
\newcommand{\CriNoDisponibles}{0}
\newcommand{\CriProfundidadArbol}{4}
\newcommand{\CriSegundosPrueba}{10.5}
\newcommand{\CriSeleccionAgr}{perfiles\_muestra}
\newcommand{\CriSeleccionAgrPuntaje}{0.144}
\newcommand{\CriSeleccionAgrSilueta}{0.433}
\newcommand{\CriSeleccionAgrValor}{1}
\newcommand{\CriSeleccionClf}{no\_conforme\_recepcion}
\newcommand{\CriSeleccionClfAUC}{0.702}
\newcommand{\CriSeleccionClfHabilidad}{0.405}
\newcommand{\CriSeleccionClfPuntaje}{0.405}
\newcommand{\CriSeleccionClfValor}{3}
\newcommand{\CriSeleccionReg}{carga\_diaria}
\newcommand{\CriSeleccionRegHabilidad}{0.290}
\newcommand{\CriSeleccionRegPuntaje}{0.193}
\newcommand{\CriSeleccionRegValor}{2}
\newcommand{\CriSupervisados}{10}
\newcommand{\CriSupervisadosEstables}{10}
\newcommand{\CriTablasConstruidas}{13}
\newcommand{\CriUmbralHopkins}{0.50}
\newcommand{\CriUmbralReglaSimplePct}{90}
\newcommand{\RegAlfaPoda}{3.798}
\newcommand{\RegAnalisis}{carga\_diaria}
\newcommand{\RegBloques}{4}
\newcommand{\RegCoefPrincipal}{puntos\_programados\_dia}
\newcommand{\RegCoefPrincipalValor}{8.85}
\newcommand{\RegCompletoEjecutado}{sí}
\newcommand{\RegConfiguracionFinal}{modelo = lineal, escalado = minmax, grupo = ampliado, objetivo\_log = no, seleccion = ninguna, codificacion = one\_hot}
\newcommand{\RegConfirmaAditivo}{sí}
\newcommand{\RegCorrParRedundante}{0.78}
\newcommand{\RegCorrResiduo}{0.10}
\newcommand{\RegCorrSimple}{0.596}
\newcommand{\RegCorridasCompleto}{128}
\newcommand{\RegDeltaPrincipal}{1.17}
\newcommand{\RegDeltaSegundo}{0.74}
\newcommand{\RegDiasEntrenamiento}{589}
\newcommand{\RegDiasFueraPct}{42.4}
\newcommand{\RegDiasFueraReferenciaPct}{56.0}
\newcommand{\RegDiasValidacion}{125}
\newcommand{\RegDidacticos}{5}
\newcommand{\RegDidacticosEstables}{5}
\newcommand{\RegFactorPrincipal}{modelo}
\newcommand{\RegFactorSegundo}{grupo}
\newcommand{\RegFactores}{6}
\newcommand{\RegFactoresMismoOptimo}{5}
\newcommand{\RegFilasReajuste}{714}
\newcommand{\RegHabilidadMejorDidactico}{0.290}
\newcommand{\RegHabilidadMinMejorDidactico}{0.283}
\newcommand{\RegHabilidadValidacion}{0.249}
\newcommand{\RegHabilidadValidacionInf}{0.154}
\newcommand{\RegHabilidadValidacionSup}{0.343}
\newcommand{\RegHiperparametros}{ninguno}
\newcommand{\RegHojasPodado}{12}
\newcommand{\RegHojasSinPoda}{95}
\newcommand{\RegInteracciones}{15}
\newcommand{\RegInteraccionesLibres}{9}
\newcommand{\RegInterceptoSimple}{27.92}
\newcommand{\RegKSeleccion}{8}
\newcommand{\RegMAEOperacion}{11.0}
\newcommand{\RegMAEReferencia}{15.4}
\newcommand{\RegMAEValidacion}{11.03}
\newcommand{\RegMaximoEntrenamiento}{135}
\newcommand{\RegMediaEntrenamiento}{38.3}
\newcommand{\RegMediaValidacion}{54.0}
\newcommand{\RegMejorCorrida}{8}
\newcommand{\RegMejorDidactico}{Lineal múltiple}
\newcommand{\RegMejorReferencia}{base\_media\_5}
\newcommand{\RegMinimoEntrenamiento}{3}
\newcommand{\RegModeloFinal}{Regresión lineal}
\newcommand{\RegModeloOtro}{Perceptrón multicapa}
\newcommand{\RegMomento}{programacion}
\newcommand{\RegObjetivo}{muestras\_dia}
\newcommand{\RegOptimo}{modelo = lineal, escalado = minmax, grupo = ampliado, objetivo\_log = no, seleccion = 8, codificacion = one\_hot}
\newcommand{\RegOrigenConfiguracion}{corrida 8 del L16}
\newcommand{\RegParRedundante}{media\_movil\_5 y media\_movil\_20}
\newcommand{\RegPendienteSimple}{1.114}
\newcommand{\RegPorEmpate}{no}
\newcommand{\RegPuestoConfirmacion}{12}
\newcommand{\RegRDosSimple}{0.355}
\newcommand{\RegRDosValidacion}{0.367}
\newcommand{\RegRMSEConfirmacion}{13.61}
\newcommand{\RegRMSEDummy}{19.40}
\newcommand{\RegRMSELinealDidactico}{13.31}
\newcommand{\RegRMSEMLPDidactico}{14.31}
\newcommand{\RegRMSEMejorCorrida}{13.38}
\newcommand{\RegRMSEMejorDidactico}{13.31}
\newcommand{\RegRMSEPodado}{16.20}
\newcommand{\RegRMSEReferencia}{18.76}
\newcommand{\RegRMSEReferenciaMax}{22.51}
\newcommand{\RegRMSEReferenciaMin}{15.67}
\newcommand{\RegRMSESinPoda}{17.33}
\newcommand{\RegRMSEValidacion}{14.59}
\newcommand{\RegRMSEValidacionInf}{12.23}
\newcommand{\RegRMSEValidacionOtro}{15.51}
\newcommand{\RegRMSEValidacionSup}{17.21}
\newcommand{\RegRejillaLinealRMSE}{13.38}
\newcommand{\RegRejillaMLPRMSE}{14.51}
\newcommand{\RegResiduoMedioValidacion}{1.85}
\newcommand{\RegRhoCompleto}{0.94}
\newcommand{\RegSNMejorCompleto}{$-$22.57}
\newcommand{\RegSNMejorCorrida}{$-$22.63}
\newcommand{\RegSNObservada}{$-$22.78}
\newcommand{\RegSNPredicha}{$-$22.76}
\newcommand{\RegSegundosCompleto}{10.3}
\newcommand{\RegSegundosLSixteen}{1.1}
\newcommand{\RegToleranciaEmpatePct}{1}
\newcommand{\RegToleranciaOperativa}{10.8}
\newcommand{\RegToleranciaOperativaPct}{20}
\newcommand{\RegToleranciaSN}{0.74}
\newcommand{\RegVariableSimple}{puntos\_programados\_dia}
\newcommand{\RegVariablesModelo}{25}

% --- 04_clasificacion_bayes ---
\newcommand{\ClfAUCCalibrado}{0.721}
\newcommand{\ClfAUCConfirmacion}{0.699}
\newcommand{\ClfAUCDesvMejorDidactico}{0.021}
\newcommand{\ClfAUCEquivalentePredicha}{0.725}
\newcommand{\ClfAUCLogisticaDidactico}{0.702}
\newcommand{\ClfAUCMLPDidactico}{0.692}
\newcommand{\ClfAUCMejorCorrida}{0.695}
\newcommand{\ClfAUCMejorDidactico}{0.702}
\newcommand{\ClfAUCPeorDidactico}{0.634}
\newcommand{\ClfAUCPodado}{0.639}
\newcommand{\ClfAUCRejillaFinal}{0.699}
\newcommand{\ClfAUCSinPoda}{0.610}
\newcommand{\ClfAUCValidacion}{0.721}
\newcommand{\ClfAUCValidacionInf}{0.697}
\newcommand{\ClfAUCValidacionOtro}{0.697}
\newcommand{\ClfAUCValidacionSinCalibrar}{0.724}
\newcommand{\ClfAUCValidacionSup}{0.745}
\newcommand{\ClfAhorroPct}{19}
\newcommand{\ClfAlertasSemana}{31.2}
\newcommand{\ClfAlertasValidacion}{843}
\newcommand{\ClfAlfaPoda}{0.000122}
\newcommand{\ClfAnalisis}{no\_conforme\_recepcion}
\newcommand{\ClfAudAUCVPPEnsayo}{0.559}
\newcommand{\ClfAudAUCVPPx}{0.779}
\newcommand{\ClfAudBrierVPPEnsayo}{0.1742}
\newcommand{\ClfAudBrierVPPx}{0.1449}
\newcommand{\ClfAudEnsayosDentroIC}{5}
\newcommand{\ClfAudFalsoRealReensayados}{0.241}
\newcommand{\ClfAudSesgoEspec}{0.0033}
\newcommand{\ClfAudVPPReal}{0.756}
\newcommand{\ClfBrierCalibrado}{0.0752}
\newcommand{\ClfBrierMejorDidactico}{0.0754}
\newcommand{\ClfBrierSinCalibrar}{0.2269}
\newcommand{\ClfCompletoEjecutado}{no}
\newcommand{\ClfConfiguracionFinal}{modelo = logistica, escalado = estandar, grupo = ampliado, desbalance = sobremuestreo, seleccion = 26, codificacion = ordinal}
\newcommand{\ClfConfirmaAditivo}{sí}
\newcommand{\ClfControlesDesarrollo}{20552}
\newcommand{\ClfCostoAlertarTodo}{5964}
\newcommand{\ClfCostoNoAlertar}{2925}
\newcommand{\ClfDeltaPrincipal}{1.030}
\newcommand{\ClfDeltaSegundo}{0.292}
\newcommand{\ClfDummyAUC}{0.500}
\newcommand{\ClfDummyBrier}{0.0813}
\newcommand{\ClfDummyPrecisionMedia}{0.089}
\newcommand{\ClfEnsayoEjemplo}{mesófilos aerobios}
\newcommand{\ClfEnsayos}{7}
\newcommand{\ClfEspecMax}{0.995}
\newcommand{\ClfEspecMin}{0.976}
\newcommand{\ClfEspecMinEnsayo}{mohos y levaduras}
\newcommand{\ClfFN}{353}
\newcommand{\ClfFP}{611}
\newcommand{\ClfFUnoUmbral}{0.325}
\newcommand{\ClfFactorPrincipal}{modelo}
\newcommand{\ClfFactorSegundo}{grupo}
\newcommand{\ClfFalsoRiesgoAlto}{0.026}
\newcommand{\ClfFalsoRiesgoBajo}{0.378}
\newcommand{\ClfFalsosPorCien}{22}
\newcommand{\ClfFilasReajuste}{29012}
\newcommand{\ClfHiperparametros}{C = 10}
\newcommand{\ClfHojaRiesgoFilas}{22}
\newcommand{\ClfHojaRiesgoFraccion}{0.773}
\newcommand{\ClfHojaRiesgoVeces}{9.0}
\newcommand{\ClfHojasPodado}{52}
\newcommand{\ClfHojasSinPoda}{546}
\newcommand{\ClfInteraccionesLibres}{9}
\newcommand{\ClfKSeleccion}{26}
\newcommand{\ClfLogisticaParametros}{136}
\newcommand{\ClfMLPActivacion}{relu}
\newcommand{\ClfMLPAlfa}{0.0010}
\newcommand{\ClfMLPCapasOcultas}{1}
\newcommand{\ClfMLPEntradas}{135}
\newcommand{\ClfMLPEpocas}{32}
\newcommand{\ClfMLPExactitudMejor}{0.921}
\newcommand{\ClfMLPFraccionValidacionPct}{15}
\newcommand{\ClfMLPMejorEpoca}{23}
\newcommand{\ClfMLPOcultas}{32}
\newcommand{\ClfMLPOptimizador}{adam}
\newcommand{\ClfMLPPaciencia}{8}
\newcommand{\ClfMLPParametros}{4385}
\newcommand{\ClfMLPPerdidaFinal}{0.222}
\newcommand{\ClfMLPTopeEpocas}{150}
\newcommand{\ClfMLPVecesParametros}{32}
\newcommand{\ClfMejorCorrida}{8}
\newcommand{\ClfMejorDidactico}{Regresión logística}
\newcommand{\ClfModeloFinal}{Regresión logística}
\newcommand{\ClfModeloMejorNivel}{regresión logística}
\newcommand{\ClfModeloOtro}{Árbol de decisión}
\newcommand{\ClfModeloPeorNivel}{perceptrón multicapa}
\newcommand{\ClfMomento}{recepcion}
\newcommand{\ClfMuestrasEntrenamiento}{22463}
\newcommand{\ClfMuestrasValidacion}{6549}
\newcommand{\ClfObjetivo}{no\_conforme}
\newcommand{\ClfOddsBaja}{0.36}
\newcommand{\ClfOddsSube}{3.65}
\newcommand{\ClfOptimo}{modelo = logistica, escalado = estandar, grupo = ampliado, desbalance = sobremuestreo, seleccion = 26, codificacion = ordinal}
\newcommand{\ClfOrigenConfiguracion}{corrida de confirmación}
\newcommand{\ClfPeorDidactico}{KNN}
\newcommand{\ClfPorEmpate}{no}
\newcommand{\ClfPositivosEntrenamiento}{1920}
\newcommand{\ClfPositivosValidacion}{585}
\newcommand{\ClfPositivosValidacionBayes}{585}
\newcommand{\ClfPosteriorMasMasMax}{0.999}
\newcommand{\ClfPosteriorMasMasMin}{0.988}
\newcommand{\ClfPosteriorMasMenosMax}{0.474}
\newcommand{\ClfPosteriorMasMenosMin}{0.059}
\newcommand{\ClfPrecisionMediaMejorDidactico}{0.254}
\newcommand{\ClfPrecisionMediaValidacion}{0.245}
\newcommand{\ClfPrecisionUmbral}{0.275}
\newcommand{\ClfPrevalenciaEntrenamientoPct}{8.5}
\newcommand{\ClfPrevalenciaValidacionPct}{8.9}
\newcommand{\ClfProbMediaCalibrado}{0.099}
\newcommand{\ClfProbMediaSinCalibrar}{0.464}
\newcommand{\ClfRecallUmbral}{0.397}
\newcommand{\ClfRecorteCeroPct}{0.0}
\newcommand{\ClfReensayados}{2186}
\newcommand{\ClfReensayoConformeEsperadoPct}{25.0}
\newcommand{\ClfReensayoConformeObservadoPct}{27.2}
\newcommand{\ClfReglaBorde}{0.22}
\newcommand{\ClfReglaCapturaPct}{51}
\newcommand{\ClfReglaGrupoPct}{17}
\newcommand{\ClfReglaTasaAltoPct}{15}
\newcommand{\ClfReglaTasaBajoPct}{78}
\newcommand{\ClfReglaVecesLimite}{1.66}
\newcommand{\ClfRejillaArbolAUC}{0.654}
\newcommand{\ClfRejillaLogisticaAUC}{0.699}
\newcommand{\ClfRevisadasPorDetectado}{2.6}
\newcommand{\ClfSNMejorCorrida}{$-$3.17}
\newcommand{\ClfSNModeloMejor}{$-$3.272}
\newcommand{\ClfSNModeloPeor}{$-$4.302}
\newcommand{\ClfSNObservada}{$-$3.117}
\newcommand{\ClfSNPredicha}{$-$2.798}
\newcommand{\ClfSegundosCompletoEstimados}{224}
\newcommand{\ClfSegundosLSixteen}{28.0}
\newcommand{\ClfSemanasValidacion}{27}
\newcommand{\ClfSensMax}{0.966}
\newcommand{\ClfSensMin}{0.897}
\newcommand{\ClfSensMinEnsayo}{Salmonella spp.}
\newcommand{\ClfSinRazonConformes}{0.338}
\newcommand{\ClfSinRazonReensayados}{346}
\newcommand{\ClfSobreDeltaAUC}{$-$0.080}
\newcommand{\ClfSobreDeltaAUCMenor}{$-$0.003}
\newcommand{\ClfSobreMasPierde}{Perceptrón multicapa}
\newcommand{\ClfSobreMenosPierde}{Regresión logística}
\newcommand{\ClfSobremuestrea}{sí}
\newcommand{\ClfSorteos}{4000}
\newcommand{\ClfTN}{5353}
\newcommand{\ClfTP}{232}
\newcommand{\ClfTasaMayoritaria}{0.915}
\newcommand{\ClfTasaTramos}{0.089}
\newcommand{\ClfToleranciaSN}{0.683}
\newcommand{\ClfUmbral}{0.167}
\newcommand{\ClfUmbralCosto}{2376}
\newcommand{\ClfUmbralCostoInf}{2196}
\newcommand{\ClfUmbralCostoSup}{2542}
\newcommand{\ClfUmbralTeorico}{0.167}
\newcommand{\ClfUmbralTeoricoCosto}{2376}
\newcommand{\ClfUsaCalibrado}{sí}
\newcommand{\ClfVPPGlobal}{0.784}
\newcommand{\ClfVPPMejor}{0.897}
\newcommand{\ClfVPPMejorEnsayo}{Salmonella spp.}
\newcommand{\ClfVPPPeor}{0.627}
\newcommand{\ClfVPPPeorEnsayo}{esterilidad}
\newcommand{\ClfVPPPeorInf}{0.428}
\newcommand{\ClfVPPPeorSup}{0.759}
\newcommand{\ClfVPPxMax}{0.995}
\newcommand{\ClfVPPxMin}{0.040}
\newcommand{\ClfVPPxPDiez}{0.597}
\newcommand{\ClfVPPxPNoventa}{0.968}
\newcommand{\ClfVariableBaja}{ensayo\_Salmonella spp.}
\newcommand{\ClfVariableSube}{tipo\_producto\_lácteo}

% --- 05_agrupamiento ---
\newcommand{\AgrAudARIDbscan}{0.007}
\newcommand{\AgrAudARIKmeans}{0.190}
\newcommand{\AgrAudClientes}{186}
\newcommand{\AgrAudPureza}{0.538}
\newcommand{\AgrAudSegmentosLatentes}{4}
\newcommand{\AgrCompARIKmDb}{0.051}
\newcommand{\AgrCompARIKmDbSinRuido}{0.000}
\newcommand{\AgrCompARISinEscalar}{0.000}
\newcommand{\AgrCompConLog}{1}
\newcommand{\AgrCompDBCombinaciones}{16}
\newcommand{\AgrCompDBConGrupos}{0}
\newcommand{\AgrCompDBCriterio}{respaldo}
\newcommand{\AgrCompDBEps}{3.42}
\newcommand{\AgrCompDBGrupos}{1}
\newcommand{\AgrCompDBMinMuestras}{12}
\newcommand{\AgrCompDBRuidoPct}{7.5}
\newcommand{\AgrCompDBSilueta}{---}
\newcommand{\AgrCompDominante}{antigüedad}
\newcommand{\AgrCompDominantePct}{100.0}
\newcommand{\AgrCompEstabilidadMax}{0.948}
\newcommand{\AgrCompEstabilidadMin}{0.934}
\newcommand{\AgrCompFilas}{186}
\newcommand{\AgrCompGrupos}{Más no conformidad y más entregas tardías, Más muestreo en planta y más tamaño de lote, Menos volumen mensual y menos muestreo en planta}
\newcommand{\AgrCompHopkins}{0.710}
\newcommand{\AgrCompK}{3}
\newcommand{\AgrCompPCAVarianzaPct}{40.2}
\newcommand{\AgrCompSegmentoMayor}{Más muestreo en planta y más tamaño de lote}
\newcommand{\AgrCompSegmentoMayorClientesPct}{39.2}
\newcommand{\AgrCompSegmentoMayorMuestrasPct}{73.6}
\newcommand{\AgrCompSilueta}{0.186}
\newcommand{\AgrCompSiluetaMax}{0.186}
\newcommand{\AgrCompSiluetaMin}{0.124}
\newcommand{\AgrCompVariables}{12}
\newcommand{\AgrComplementario}{segmentos\_clientes}
\newcommand{\AgrComplementarioSupera}{sí}
\newcommand{\AgrEfectoFilas}{0.001}
\newcommand{\AgrEfectoLog}{$-$0.286}
\newcommand{\AgrKMax}{10}
\newcommand{\AgrKMin}{2}
\newcommand{\AgrPuntajeCompEtapa}{0.124}
\newcommand{\AgrPuntajeSelEtapa}{0.049}
\newcommand{\AgrRuidoMaximoPct}{15}
\newcommand{\AgrRuidoRespaldoPct}{10}
\newcommand{\AgrSelARIKmDb}{0.135}
\newcommand{\AgrSelARIKmDbSinRuido}{0.064}
\newcommand{\AgrSelARISinEscalar}{0.000}
\newcommand{\AgrSelConLog}{5}
\newcommand{\AgrSelDBCombinaciones}{16}
\newcommand{\AgrSelDBConGrupos}{10}
\newcommand{\AgrSelDBCriterio}{regla}
\newcommand{\AgrSelDBEps}{1.59}
\newcommand{\AgrSelDBGrupos}{2}
\newcommand{\AgrSelDBMinMuestras}{16}
\newcommand{\AgrSelDBRuidoPct}{12.7}
\newcommand{\AgrSelDBSilueta}{0.314}
\newcommand{\AgrSelDominante}{volumen de muestra}
\newcommand{\AgrSelDominantePct}{95.0}
\newcommand{\AgrSelEstabilidadMax}{0.996}
\newcommand{\AgrSelEstabilidadMin}{0.981}
\newcommand{\AgrSelFilas}{5000}
\newcommand{\AgrSelGrupos}{Recepción tardía y transporte más frío, Transporte largo, Recepción temprana, Plazo prometido largo y entrega lenta}
\newcommand{\AgrSelHopkins}{0.714}
\newcommand{\AgrSelK}{4}
\newcommand{\AgrSelPCAVarianzaPct}{38.1}
\newcommand{\AgrSelSilueta}{0.148}
\newcommand{\AgrSelSiluetaMax}{0.148}
\newcommand{\AgrSelSiluetaMin}{0.120}
\newcommand{\AgrSelVariables}{8}
\newcommand{\AgrSeleccionado}{perfiles\_muestra}
\newcommand{\AgrSemillasEstabilidad}{5}
\newcommand{\AgrSiluetaCribado}{0.433}
\newcommand{\AgrSiluetaSinLog}{0.434}
\newcommand{\AgrToleranciaSilueta}{0.02}
\newcommand{\AgrUmbralNombre}{0.5}

% --- 06_evaluacion_final ---
\newcommand{\FinAUCPrueba}{0.724}
\newcommand{\FinAUCPruebaInf}{0.705}
\newcommand{\FinAUCPruebaSup}{0.743}
\newcommand{\FinAgrHopkins}{0.714}
\newcommand{\FinAgrK}{4}
\newcommand{\FinAgrSegmentos}{3}
\newcommand{\FinAgrSilueta}{0.148}
\newcommand{\FinAgrSiluetaCribado}{0.433}
\newcommand{\FinAhorroPct}{15}
\newcommand{\FinAlertasSemana}{21.2}
\newcommand{\FinAudBrierVPPEnsayo}{0.1586}
\newcommand{\FinAudBrierVPPx}{0.1381}
\newcommand{\FinAudVPPReal}{0.800}
\newcommand{\FinBrierPrueba}{0.0753}
\newcommand{\FinCambioRMSEPct}{$-$15}
\newcommand{\FinClfModelo}{Regresión logística}
\newcommand{\FinCostoAlertarTodo}{7882}
\newcommand{\FinCostoNoAlertar}{3875}
\newcommand{\FinCostoPorMuestraPrueba}{0.381}
\newcommand{\FinCostoPrueba}{3295}
\newcommand{\FinCostoRetro}{3231}
\newcommand{\FinDiasFinDeAnio}{5}
\newcommand{\FinDiasFueraPct}{37.9}
\newcommand{\FinDiasFueraReferenciaPct}{57.1}
\newcommand{\FinDiasPrueba}{203}
\newcommand{\FinEnsayosPocosPositivos}{2}
\newcommand{\FinFN}{523}
\newcommand{\FinFP}{680}
\newcommand{\FinFUnoPrueba}{0.295}
\newcommand{\FinHabilidadPrueba}{0.334}
\newcommand{\FinHabilidadPruebaInf}{0.247}
\newcommand{\FinHabilidadPruebaSup}{0.408}
\newcommand{\FinMAEPrueba}{9.80}
\newcommand{\FinMediaPrueba}{44.1}
\newcommand{\FinMinimoPositivos}{30}
\newcommand{\FinMuestrasPrueba}{8657}
\newcommand{\FinPositivosPrueba}{775}
\newcommand{\FinPositivosReensayados}{691}
\newcommand{\FinPrecisionMediaPrueba}{0.243}
\newcommand{\FinPrecisionPrueba}{0.270}
\newcommand{\FinPrecisionPruebaInf}{0.242}
\newcommand{\FinPrecisionPruebaSup}{0.299}
\newcommand{\FinPrevalenciaPrueba}{0.090}
\newcommand{\FinProbMediaPrueba}{0.090}
\newcommand{\FinPruebaDesde}{2026-01-02}
\newcommand{\FinPruebaHasta}{2026-10-30}
\newcommand{\FinRDosPrueba}{0.529}
\newcommand{\FinRMSEPrueba}{12.47}
\newcommand{\FinRMSEPruebaInf}{11.20}
\newcommand{\FinRMSEPruebaSup}{13.66}
\newcommand{\FinRecallPrueba}{0.325}
\newcommand{\FinRecallPruebaInf}{0.293}
\newcommand{\FinRecallPruebaSup}{0.357}
\newcommand{\FinReensayoConformeEsperadoPct}{24.1}
\newcommand{\FinReensayoConformeObservadoInfPct}{20.3}
\newcommand{\FinReensayoConformeObservadoPct}{23.3}
\newcommand{\FinReensayoConformeObservadoSupPct}{26.6}
\newcommand{\FinReensayosLecturaPrueba}{698}
\newcommand{\FinReensayosRecibidosAntes}{7}
\newcommand{\FinRegModelo}{Regresión lineal}
\newcommand{\FinRegReferencia}{base\_media\_5}
\newcommand{\FinReglaBajoBorde}{93}
\newcommand{\FinReglaCapturaPct}{55}
\newcommand{\FinReglaReensayosConRecuento}{593}
\newcommand{\FinReglaTasaAltoPct}{12}
\newcommand{\FinReglaTasaBajoInfPct}{68}
\newcommand{\FinReglaTasaBajoPct}{77}
\newcommand{\FinReglaTasaBajoSupPct}{85}
\newcommand{\FinResiduoFinDeAnio}{$-$12.5}
\newcommand{\FinResiduoResto}{0.6}
\newcommand{\FinSemanasPrueba}{44}
\newcommand{\FinTN}{7202}
\newcommand{\FinTP}{252}
\newcommand{\FinUmbral}{0.167}
\newcommand{\FinUmbralRetro}{0.140}
\newcommand{\FinUnoDeCadaFalso}{5}
\newcommand{\FinVPPxMedio}{0.794}

 define las macros;
% cada tabla se incluye con \input{generado/tabla_<nombre>}.
%
% ===== 01_datos_exploracion: 91 macros =====
% \DatBloquesCV                      = 4            bloques temporales de la validación cruzada
% \DatClientes                       = 220          clientes en la tabla de clientes
% \DatClientesConMuestras            = 219          clientes con al menos una muestra
% \DatColumnasContrato               = 65           columnas declaradas en el contrato de datos
% \DatColumnasObligatorias           = 15           columnas obligatorias del contrato
% \DatColumnasOpcionales             = 50           columnas opcionales del contrato
% \DatControles                      = 26236        controles de calidad (positivos y negativos)
% \DatCortePrueba                    = 2026-01-01   fecha de corte entre validación y prueba
% \DatCorteValidacion                = 2025-07-01   fecha de corte entre entrenamiento y validación
% \DatCostoFN                        = 5            costo de un falso negativo (no conforme sin alerta)
% \DatCostoFP                        = 1            costo de un falso positivo (alerta innecesaria)
% \DatDesde                          = 2023-01-02   primera fecha de recepción
% \DatDiasCalendario                 = 1398         días de la tabla de calendario
% \DatDuplicadas                     = 148          filas crudas de muestras duplicadas exactas
% \DatFaltanOpcionales               = 0            columnas opcionales que no trae el origen
% \DatFechaFin                       = 2026-10-30   fin del periodo de datos
% \DatFechaInicio                    = 2023-01-02   inicio del periodo de datos
% \DatHasta                          = 2026-10-30   última fecha de recepción
% \DatHuella                         = f82dd7cc9bcb huella SHA-256 abreviada de las tablas crudas
% \DatMomentos                       = 9            momentos del proceso declarados en el contrato
% \DatMuestrasCrudas                 = 38994        filas crudas de la tabla de muestras
% \DatOrigen                         = simulacion   origen de los datos (simulacion o real)
% \DatRecuentosNoNumericos           = 28716        recuentos crudos que no son un número simple
% \DatRecuentosNoNumericosPct        = 75.9         porcentaje de recuentos no numéricos
% \DatSemilla                        = 42           semilla del proyecto
% \DatTablasContrato                 = 4            tablas del contrato de datos
% \DatValoresImposiblesCrudos        = 366          valores fuera de su rango físico en las filas crudas, con las duplicadas incluidas
% \DatValoresImposiblesEnDuplicadas  = 1            de ellos, los que caen en filas duplicadas (la auditoría de limpieza cuenta tras quitarlas)
% \DatVisitas                        = 1812         visitas de muestreo en planta
% \ExpCVSemanal                      = 0.30         coeficiente de variación del volumen semanal
% \ExpCategoriaMasNiveles            = codigo\_producto categórica con más niveles
% \ExpClientesDesarrollo             = 203          clientes con muestras en el desarrollo
% \ExpClientesMitad                  = 28           clientes más grandes que aportan la mitad del volumen
% \ExpClientesMitadPct               = 13.8         esos clientes como porcentaje del total
% \ExpColumnasConFaltantes           = 25           columnas con algún faltante en el desarrollo
% \ExpColumnasIntegradas             = 69           columnas de la tabla integrada por muestra
% \ExpColumnasSoloEstructural        = 17           columnas cuyos faltantes son todos estructurales
% \ExpConformesPorNoConforme         = 11           conformes por cada no conforme (desarrollo)
% \ExpDecimoClientesPct              = 41.1         volumen (%) del 10 % de clientes más grande
% \ExpDiferenciaEstandarMax          = 0.34         diferencia de medias en desviaciones estándar
% \ExpDuplicadosEliminados           = 148          muestras duplicadas eliminadas
% \ExpEnsayoMasTardio                = mohos y levaduras ensayo con más entregas tardías
% \ExpEnsayoMasTardioPct             = 16.6         entregas tardías (%) de ese ensayo
% \ExpEnsayoMenosTardio              = Staphylococcus aureus ensayo con menos entregas tardías
% \ExpEnsayoMenosTardioPct           = 8.8          entregas tardías (%) de ese ensayo
% \ExpEntregaMediana                 = 6            mediana de días de entrega
% \ExpEntregaPNoventa                = 12           percentil 90 de días de entrega
% \ExpEntregaTardiaPct               = 12.6         porcentaje de informes entregados tarde
% \ExpFilasEntrenamiento             = 23151        muestras del bloque de entrenamiento
% \ExpFilasPrueba                    = 8947         muestras del bloque de prueba (solo se cuentan)
% \ExpFilasValidacion                = 6748         muestras del bloque de validación
% \ExpHorasCuartoAltoPct             = 12.6         no conformes (%) en el cuarto con mayor duración del transporte
% \ExpHorasCuartoBajoPct             = 7.0          no conformes (%) en el cuarto con menor duración del transporte
% \ExpIndiceMesAlto                  = 1.26         índice estacional del mes más cargado
% \ExpIndiceMesBajo                  = 0.83         índice estacional del mes más liviano
% \ExpMedianaMuestrasCliente         = 73           mediana de muestras por cliente
% \ExpMesAlto                        = nov          mes más cargado
% \ExpMesBajo                        = ene          mes más liviano
% \ExpMuestrasCrudas                 = 38994        muestras crudas leídas por la limpieza
% \ExpMuestrasDesarrollo             = 29899        muestras del periodo de desarrollo
% \ExpMuestrasLimpias                = 38846        muestras tras la limpieza
% \ExpMuestrasSemana                 = 192          muestras por semana en promedio (semanas completas)
% \ExpMuestrasSemanaMax              = 366          semana completa con más muestras
% \ExpMuestrasSemanaMin              = 61           semana completa con menos muestras
% \ExpNivelesCategoria               = 1169         niveles de esa categórica
% \ExpParCorrelacionado              = dias\_prometidos y limite\_ufc par de variables de entrada más correlacionado
% \ExpPeorFaltanteColumna            = actividad\_agua columna con más faltantes no estructurales
% \ExpPeorFaltantePct                = 9.8          porcentaje de faltantes no estructurales de esa columna
% \ExpPrimerAnio                     = 2023         primer año del desarrollo
% \ExpProblemaMayor                  = recuento censurado (``<N'' o ``>N'') tipo de problema de limpieza que más filas toca
% \ExpProblemaMayorFilas             = 18816        filas que toca el problema de limpieza mayor
% \ExpReglasConHallazgos             = 18           reglas de limpieza con hallazgos (sin faltantes)
% \ExpRhoObjetivo                    = 0.07         su correlación de Spearman con no conforme
% \ExpRhoPar                         = $-$0.28      correlación de Spearman de ese par
% \ExpSectorMayor                    = cosmético    línea de negocio con más volumen
% \ExpSectorMayorPct                 = 36.3         porcentaje de muestras de esa línea
% \ExpSectorMenor                    = central de mezclas línea de negocio con menos volumen
% \ExpSectorMenorPct                 = 8.6          porcentaje de muestras de esa línea
% \ExpSemanas                        = 157          semanas del periodo de desarrollo
% \ExpTablasLimpias                  = 5            tablas limpias guardadas
% \ExpTasaMensualMaxPct              = 12.7         tasa mensual máxima de no conformes (%)
% \ExpTasaMensualMinPct              = 5.5          tasa mensual mínima de no conformes (%)
% \ExpTasaNoConformePct              = 8.6          porcentaje de no conformes entre aceptadas (desarrollo)
% \ExpTasaPrimerAnioPct              = 7.9          no conformes (%) en el primer año
% \ExpTasaRechazoPct                 = 3.0          porcentaje de muestras rechazadas en recepción (desarrollo)
% \ExpTasaUltimoAnioPct              = 8.7          no conformes (%) en el último año
% \ExpTemperaturaCuartoAltoPct       = 12.8         no conformes (%) en el cuarto con mayor temperatura de transporte
% \ExpTemperaturaCuartoBajoPct       = 6.9          no conformes (%) en el cuarto con menor temperatura de transporte
% \ExpUltimoAnio                     = 2025         último año del desarrollo
% \ExpVariableCorrObjetivo           = actividad\_agua variable más correlacionada con no conforme
% \ExpVariableMasSepara              = horas\_transporte variable de llegada que más separa no conformes de conformes
% --- tablas de 01_datos_exploracion ---
% tabla_auditoria_limpieza.tex (7 filas): Correcciones de la limpieza por tipo de problema
% tabla_cobertura_opcionales.tex (4 filas): Columnas opcionales que trae el origen, por tabla
% tabla_columnas_por_momento.tex (9 filas): Columnas del contrato que se conocen en cada momento
% tabla_diccionario_datos.tex (65 filas): Diccionario de datos completo (contrato)
% tabla_filas_crudas.tex (5 filas): Filas de cada tabla cruda
% tabla_particion.tex (3 filas): Muestras por bloque temporal
%
% ===== 02_base_real: 145 macros =====
% \RealArbolGiniAntes            = 0.0690       Gini de la raíz del árbol
% \RealArbolGiniDespues          = 0.0666       Gini ponderado tras el primer corte
% \RealArbolGiniReduccion        = 0.0023       reducción del Gini en el primer corte
% \RealArbolPrimerCorte          = tasa\_cliente > 0.0735105 primer corte del árbol
% \RealCalcIntercepto            = $-$3.677     intercepto de la logística
% \RealCalcLogit                 = 2.047        logit de la muestra del cálculo ilustrativo
% \RealCalcMayor                 = tasa\_cliente variable con el mayor aporte en el cálculo
% \RealCalcMayorAporte           = 5.442        su aporte al logit
% \RealCalcProba                 = 0.886        probabilidad de la muestra del cálculo ilustrativo
% \RealCalcTerminos              = 9            variables activas en el cálculo ilustrativo
% \RealCampanaMuestrasMediana    = 14           muestras por campaña (mediana)
% \RealCampanaUnDiaPct           = 69           campañas de un solo día (%)
% \RealCampanas                  = 385          campañas de muestreo en planta
% \RealClientes                  = 82           clientes con muestras en el desarrollo
% \RealClientesDispersion        = 28           clientes con lotes suficientes (todos los tipos)
% \RealClientesPlanta            = 8            clientes con lotes de planta suficientes
% \RealClientesProducto          = 22           clientes con lotes de producto suficientes
% \RealCompuestoPlanta           = 59.98        índice compuesto en planta
% \RealCompuestoProducto         = 6.00         índice compuesto en producto
% \RealCongeladoHuella           = 10900224adf7 primeros caracteres de la huella del modelo
% \RealCongeladoModelo           = Regresión logística (C = 0.01) modelo congelado
% \RealCostoFN                   = 5            costo supuesto de un falso negativo
% \RealCostoFP                   = 1            costo supuesto de un falso positivo
% \RealDesde                     = 2024-06-04   primer ingreso del desarrollo
% \RealDispLotesPlanta           = 1.20         índice de lotes de planta
% \RealDispLotesProducto         = 0.73         índice de lotes de producto
% \RealDispLotesTodos            = 0.90         mediana del índice de lotes por semana
% \RealDispMuestrasMax           = 407          máximo índice de muestras por semana
% \RealDispMuestrasPlanta        = 71.16        índice de muestras de planta
% \RealDispMuestrasProducto      = 4.25         índice de muestras de producto
% \RealDispMuestrasTodos         = 7.23         mediana del índice de muestras por semana
% \RealFilas                     = 120114       filas de desarrollo de la exportación real
% \RealHasta                     = 2025-12-19   último ingreso del desarrollo
% \RealHojaFilas                 = 51           muestras en la hoja de mayor riesgo
% \RealHojaFraccion              = 0.314        fracción de no conformes en esa hoja
% \RealHojaVeces                 = 8.8          veces la tasa general
% \RealLoteDeltaGeoPlanta        = 448          AIC de la geométrica sobre el mejor (planta)
% \RealLoteDeltaGeoProducto      = 5.4          AIC de la geométrica sobre el mejor (producto)
% \RealLoteMaxPlanta             = 675          lote de planta más grande
% \RealLoteMaxProducto           = 60           lote de producto más grande
% \RealLoteMediaPlanta           = 18.7         tamaño medio de los lotes de planta
% \RealLoteMediaProducto         = 4.1          tamaño medio de los lotes de producto
% \RealLoteMedianaPlanta         = 6            tamaño mediano de los lotes de planta
% \RealLoteMedianaProducto       = 3            tamaño mediano de los lotes de producto
% \RealLoteMejorPlanta           = 1 + binomial negativa mejor modelo del tamaño de lote en planta
% \RealLoteMejorProducto         = 1 + binomial negativa mejor modelo del tamaño de lote en producto
% \RealLoteRPlanta               = 0.39         r de la binomial negativa en planta
% \RealLoteRProducto             = 0.88         r de la binomial negativa en producto
% \RealLotes                     = 2543         lotes de desarrollo (cliente, día y tipo)
% \RealLotesPlanta               = 762          lotes de planta
% \RealLotesProducto             = 1781         lotes de producto
% \RealMinimoLotes               = 20           lotes mínimos de un cliente para medir su dispersión
% \RealMuestras                  = 21549        muestras de desarrollo de la exportación real
% \RealNivel                     = 0.05         nivel de la prueba de dispersión
% \RealNoConformePct             = 3.7          no conformidad de las muestras de desarrollo (%)
% \RealOddsBaja                  = previas\_lote variable que más baja el riesgo
% \RealOddsBajaValor             = 0.70         su razón de odds
% \RealOddsSube                  = tasa\_cliente variable que más sube el riesgo
% \RealOddsSubeValor             = 1.76         su razón de odds
% \RealPlantaPct                 = 66.2         muestras de muestreo en planta (%)
% \RealPruebaAP                  = 0.072        PR-AUC del modelo en la prueba
% \RealPruebaAPConocidos         = 0.073        PR-AUC en clientes conocidos
% \RealPruebaAPInf               = 0.059        límite inferior del PR-AUC en la prueba
% \RealPruebaAPNuevos            = 0.097        PR-AUC en clientes nuevos
% \RealPruebaAPRegla             = 0.064        PR-AUC de la regla en la prueba
% \RealPruebaAPReglaInf          = 0.053        límite inferior del PR-AUC de la regla
% \RealPruebaAPReglaSup          = 0.081        límite superior del PR-AUC de la regla
% \RealPruebaAPSup               = 0.091        límite superior del PR-AUC en la prueba
% \RealPruebaAUC                 = 0.769        ROC-AUC del modelo en la prueba
% \RealPruebaAUCInf              = 0.743        límite inferior del ROC-AUC en la prueba
% \RealPruebaAUCSup              = 0.794        límite superior del ROC-AUC en la prueba
% \RealPruebaClientes            = 72           clientes en la prueba
% \RealPruebaClientesNuevos      = 17           clientes de la prueba sin desarrollo
% \RealPruebaCosto               = 1718         costo en la prueba con el umbral fijo
% \RealPruebaCostoNada           = 1710         costo en la prueba de no alertar nunca
% \RealPruebaFN                  = 339          falsos negativos en la prueba
% \RealPruebaFP                  = 23           falsos positivos en la prueba
% \RealPruebaFUno                = 0.016        F1 en la prueba con el umbral fijo
% \RealPruebaMuestras            = 16325        muestras con dictamen en la prueba de 2026
% \RealPruebaNuevasPct           = 29           muestras de prueba de clientes nuevos (%)
% \RealPruebaPositivos           = 342          no conformes en la prueba
% \RealPruebaPrecision           = 0.115        precisión en la prueba con el umbral fijo
% \RealPruebaPrevalencia         = 0.021        prevalencia de la prueba
% \RealPruebaPrevalenciaPct      = 2.1          no conformidad en la prueba (%)
% \RealPruebaRecall              = 0.009        recall en la prueba con el umbral fijo
% \RealPruebaSinHistoriaPct      = 2.2          muestras de prueba sin dictamen previo del cliente (%)
% \RealPruebaTN                  = 15960        verdaderos negativos en la prueba
% \RealPruebaTP                  = 3            verdaderos positivos en la prueba
% \RealPruebaVeces               = 3.4          PR-AUC sobre la prevalencia
% \RealRiesgoAPArbol             = 0.073        PR-AUC del mejor árbol en validación
% \RealRiesgoAPArbolProfundoEnt  = 0.212        PR-AUC del árbol más profundo en entrenamiento
% \RealRiesgoAPArbolProfundoVal  = 0.073        PR-AUC del árbol más profundo en validación
% \RealRiesgoAPBase              = 0.041        PR-AUC de la referencia en validación
% \RealRiesgoAPElegido           = 0.109        PR-AUC del elegido en validación
% \RealRiesgoAPHistoria          = 0.105        PR-AUC de la tasa histórica en validación
% \RealRiesgoAPLogistica         = 0.109        PR-AUC de la mejor regresión logística
% \RealRiesgoAPLogisticaEnt      = 0.114        PR-AUC de la mejor logística en entrenamiento
% \RealRiesgoAUCArbol            = 0.683        ROC-AUC del mejor árbol
% \RealRiesgoAUCElegido          = 0.689        ROC-AUC del elegido en validación
% \RealRiesgoAUCHistoria         = 0.673        ROC-AUC de la tasa histórica del cliente
% \RealRiesgoAUCLogistica        = 0.689        ROC-AUC de la mejor regresión logística
% \RealRiesgoArbol               = profundidad 5 mejor configuración del árbol
% \RealRiesgoElegido             = Regresión logística (C = 0.01) configuración elegida por PR-AUC en validación
% \RealRiesgoEntrenamiento       = 15628        muestras de entrenamiento del clasificador real
% \RealRiesgoEntrenamientoPct    = 73           entrenamiento sobre el desarrollo (%)
% \RealRiesgoHojaEsperados       = 1.8          no conformes esperados en una hoja mínima
% \RealRiesgoLogistica           = C = 0.01     mejor configuración logística
% \RealRiesgoMinimoHoja          = 50           mínimo de muestras por hoja del árbol
% \RealRiesgoPrevalenciaEntPct   = 3.6          no conformidad en entrenamiento (%)
% \RealRiesgoPrevalenciaPct      = 4.1          no conformidad en validación (%)
% \RealRiesgoSinDictamen         = 98           muestras de desarrollo sin dictamen
% \RealRiesgoValidacion          = 5823         muestras de validación del clasificador real
% \RealRiesgoValidacionPct       = 27           validación sobre el desarrollo (%)
% \RealRiesgoVariables           = 9            variables conocidas al ingreso
% \RealSegClientes               = 46           clientes en el agrupamiento real
% \RealSegK                      = 4            número de segmentos de clientes reales
% \RealSegMayor                  = Más muestras por semana y más dispersión semanal segmento real con más muestras
% \RealSegMayorMuestrasPct       = 56.3         sus muestras (%)
% \RealSegMinimoMuestras         = 20           muestras mínimas para entrar al agrupamiento
% \RealSegNombreCuatro           = Más muestras por semana y más dispersión semanal nombre del segmento real 4
% \RealSegNombreDos              = Más muestreo en planta y más tamaño de lote nombre del segmento real 2
% \RealSegNombreTres             = Menos muestreo en planta y menos lotes por semana nombre del segmento real 3
% \RealSegNombreUno              = Más semanas con envío y más lotes por semana nombre del segmento real 1
% \RealSegSilueta                = 0.374        silueta de los segmentos reales
% \RealSimIndiceAlto             = 3.7          índice de muestras de la simulación con el lote medio mayor
% \RealSimIndiceBajo             = 2.2          índice de muestras de la simulación con el lote medio menor
% \RealSobreLotesPlantaPct       = 25           clientes con lotes sobredispersos en planta (%)
% \RealSobreLotesProductoPct     = 5            clientes con lotes sobredispersos en producto (%)
% \RealSobreMuestrasPlantaPct    = 100          clientes con muestras sobredispersas en planta (%)
% \RealSobreMuestrasProductoPct  = 95           clientes con muestras sobredispersas en producto (%)
% \RealUmbral                    = 0.61         umbral elegido por costo en validación
% \RealUmbralAlertas             = 41           alertas en validación con el umbral elegido
% \RealUmbralAlertasSemana       = 1.6          alertas por semana en validación
% \RealUmbralCosto               = 1152         costo en validación con el umbral elegido
% \RealUmbralCostoNada           = 1195         costo en validación de no alertar nunca
% \RealUmbralCostoTodo           = 5584         costo en validación de alertar todas
% \RealUmbralFN                  = 225          falsos negativos en validación
% \RealUmbralFP                  = 27           falsos positivos en validación
% \RealUmbralFUno                = 0.100        F1 en validación con el umbral
% \RealUmbralPrecision           = 0.341        precisión en validación con el umbral
% \RealUmbralRecall              = 0.059        recall en validación con el umbral
% \RealUmbralTN                  = 5557         verdaderos negativos en validación
% \RealUmbralTP                  = 14           verdaderos positivos en validación
% \RealUmbralTeorico             = 0.167        umbral teórico por costo
% \RealValidacionVeces           = 2.7          PR-AUC de validación sobre su prevalencia
% --- tablas de 02_base_real ---
% tabla_real_calculo_logistica.tex (10 filas): Cálculo ilustrativo de la regresión logística
% tabla_real_configuraciones.tex (16 filas): Configuraciones del clasificador real en entrenamiento y validación
% tabla_real_dispersion.tex (2 filas): Índice de dispersión semanal por cliente y tipo de lote
% tabla_real_odds.tex (8 filas): Variables con mayor peso en la logística real
% tabla_real_particiones.tex (2 filas): Muestras de desarrollo por partición
% tabla_real_prueba.tex (3 filas): Prueba de 2026 con lo congelado
% tabla_real_reglas_arbol.tex (5 filas): Hojas de mayor riesgo del árbol real
% tabla_real_riesgo_particiones.tex (3 filas): Particiones del clasificador real
% tabla_real_segmentos.tex (4 filas): Segmentos de clientes reales
% tabla_real_simulacion.tex (4 filas): Mecanismo de llegada en la simulación y en la exportación real
% tabla_real_tamano_lote.tex (6 filas): Modelos del tamaño de lote por tipo
%
% ===== 03_cribado_regresion: 139 macros =====
% \CriAgrCandidatos               = 3            candidatos de agrupamiento cribados
% \CriAgrComplementario           = segmentos\_clientes análisis de agrupamiento complementario (fuera de la regla)
% \CriAnalisisConTabla            = 13           análisis del catálogo que llevan tabla de modelado
% \CriAuditoriasSuperadas         = 13           tablas que superan la auditoría de fugas
% \CriBayes                       = 3            análisis bayesianos (no se criban)
% \CriCandidatos                  = 16           análisis candidatos del catálogo
% \CriCandidatosAgr               = 3            candidatos de tipo agrupamiento
% \CriCandidatosBayes             = 3            candidatos de tipo bayes
% \CriCandidatosClf               = 5            candidatos de tipo clasificacion
% \CriCandidatosReg               = 5            candidatos de tipo regresion
% \CriClfMasHabil                 = falso\_positivo\_lectura clasificador con la mayor habilidad bruta en el cribado
% \CriClfMasHabilHabilidad        = 0.689        su habilidad media (2 AUC - 1)
% \CriClfMasHabilProporcion       = 95           porcentaje de su habilidad que logra la mejor variable sola
% \CriCribados                    = 13           candidatos que pasan por la prueba rápida
% \CriDescartadosReglaSimple      = 3            candidatos descartados por la regla simple
% \CriElegibles                   = 10           candidatos elegibles según la regla
% \CriElegidos                    = 3            análisis elegidos para el modelado a fondo
% \CriHopkinsMax                  = 0.86         Hopkins máximo entre candidatos de agrupamiento
% \CriHopkinsMin                  = 0.72         Hopkins mínimo entre candidatos de agrupamiento
% \CriMasInestable                = cliente\_inactivo\_trimestre candidato con la habilidad más inestable entre bloques
% \CriMasInestableMax             = 0.708        su habilidad máxima por bloque
% \CriMasInestableMin             = 0.442        su habilidad mínima por bloque
% \CriMejorSilueta                = 0.433        mejor silueta del cribado
% \CriMejorSiluetaAnalisis        = perfiles\_muestra candidato de agrupamiento con la mejor silueta
% \CriMejorSiluetaK               = 2            k de la mejor silueta del cribado
% \CriMinimoFilas                 = 300          filas mínimas de entrenamiento de un supervisado
% \CriMinimoFilasAgrupamiento     = 50           filas mínimas de un agrupamiento
% \CriMinimoPositivos             = 100          positivos mínimos de un clasificador
% \CriNoDisponibles               = 0            análisis no disponibles por columnas faltantes
% \CriProfundidadArbol            = 4            profundidad del árbol rápido del cribado
% \CriSegundosPrueba              = 9.6          segundos de la prueba rápida completa
% \CriSeleccionAgr                = perfiles\_muestra análisis de agrupamiento elegido por la regla
% \CriSeleccionAgrPuntaje         = 0.144        puntaje del agrupamiento elegido
% \CriSeleccionAgrSilueta         = 0.433        silueta del agrupamiento elegido
% \CriSeleccionAgrValor           = 1            valor de negocio del agrupamiento elegido
% \CriSeleccionClf                = no\_conforme\_recepcion análisis de clasificacion elegido por la regla
% \CriSeleccionClfAUC             = 0.702        AUC del clasificador elegido
% \CriSeleccionClfHabilidad       = 0.405        habilidad media del clasificacion elegido
% \CriSeleccionClfPuntaje         = 0.405        puntaje del clasificacion elegido
% \CriSeleccionClfValor           = 3            valor de negocio del clasificacion elegido
% \CriSeleccionReg                = carga\_diaria análisis de regresion elegido por la regla
% \CriSeleccionRegHabilidad       = 0.290        habilidad media del regresion elegido
% \CriSeleccionRegPuntaje         = 0.193        puntaje del regresion elegido
% \CriSeleccionRegValor           = 2            valor de negocio del regresion elegido
% \CriSupervisados                = 10           candidatos supervisados cribados
% \CriSupervisadosEstables        = 10           supervisados que superan a su base en todos los bloques
% \CriTablasConstruidas           = 13           tablas de modelado construidas
% \CriUmbralHopkins               = 0.50         umbral del estadístico de Hopkins
% \CriUmbralReglaSimplePct        = 90           umbral (%) de la regla simple
% \RegAlfaPoda                    = 3.798        ccp_alpha elegido para el árbol didáctico
% \RegAnalisis                    = carga\_diaria análisis de regresión modelado a fondo
% \RegBloques                     = 4            bloques temporales de la validación cruzada
% \RegCoefPrincipal               = puntos\_programados\_dia variable con el mayor coeficiente estandarizado
% \RegCoefPrincipalValor          = 8.85         valor de ese coeficiente
% \RegCompletoEjecutado           = sí           se ejecutó el factorial completo
% \RegConfiguracionFinal          = modelo = lineal, escalado = minmax, grupo = ampliado, objetivo\_log = no, seleccion = ninguna, codificacion = one\_hot configuración congelada
% \RegConfirmaAditivo             = sí           la confirmación cae dentro de la tolerancia
% \RegCorrParRedundante           = 0.78         su correlación absoluta
% \RegCorrResiduo                 = 0.10         correlación entre la predicción y el error absoluto en validación
% \RegCorrSimple                  = 0.596        correlación de Pearson de esa variable con el objetivo
% \RegCorridasCompleto            = 128          corridas del factorial completo
% \RegDeltaPrincipal              = 1.17         delta (dB) del factor más influyente
% \RegDeltaSegundo                = 0.74         delta (dB) del segundo factor
% \RegDiasEntrenamiento           = 589          días de entrenamiento
% \RegDiasFueraPct                = 42.4         días de validación fuera de tolerancia (%), modelo
% \RegDiasFueraReferenciaPct      = 56.0         días de validación fuera de tolerancia (%), referencia
% \RegDiasValidacion              = 125          días de validación
% \RegDidacticos                  = 5            modelos de la pasada didáctica
% \RegDidacticosEstables          = 5            modelos didácticos que superan a la referencia en todos los bloques
% \RegFactorPrincipal             = modelo       factor más influyente (mayor delta)
% \RegFactorSegundo               = grupo        segundo factor más influyente
% \RegFactores                    = 6            factores del L16 de regresión
% \RegFactoresMismoOptimo         = 5            factores con el mismo mejor nivel en ambos diseños
% \RegFilasReajuste               = 714          filas del reajuste final (entrenamiento + validación)
% \RegHabilidadMejorDidactico     = 0.290        su habilidad media frente a la referencia
% \RegHabilidadMinMejorDidactico  = 0.283        su habilidad mínima por bloque
% \RegHabilidadValidacion         = 0.249        habilidad del modelo final en validación
% \RegHabilidadValidacionInf      = 0.154        límite inferior del IC 95 % de la habilidad en validación
% \RegHabilidadValidacionSup      = 0.343        límite superior del IC 95 % de la habilidad en validación
% \RegHiperparametros             = ninguno      hiperparámetros congelados
% \RegHojasPodado                 = 12           hojas del árbol podado
% \RegHojasSinPoda                = 95           hojas del árbol de regresión sin poda
% \RegInteracciones               = 15           interacciones de dos factores
% \RegInteraccionesLibres         = 9            interacciones que quedan libres en el L16
% \RegInterceptoSimple            = 27.92        intercepto de la recta simple
% \RegKSeleccion                  = 8            variables que conserva la selección por F
% \RegMAEOperacion                = 11.0         error absoluto medio del modelo en validación (muestras)
% \RegMAEReferencia               = 15.4         error absoluto medio de la referencia en validación (muestras)
% \RegMAEValidacion               = 11.03        MAE del modelo final en validación
% \RegMaximoEntrenamiento         = 135          máximo diario observado (entrenamiento)
% \RegMediaEntrenamiento          = 38.3         muestras por día en promedio (entrenamiento)
% \RegMediaValidacion             = 54.0         muestras por día hábil en promedio (validación)
% \RegMejorCorrida                = 8            número de la mejor corrida observada del L16
% \RegMejorDidactico              = Lineal múltiple mejor modelo de la pasada didáctica
% \RegMejorReferencia             = base\_media\_5 mejor referencia ingenua (base de la habilidad)
% \RegMinimoEntrenamiento         = 3            mínimo diario observado (entrenamiento)
% \RegModeloFinal                 = Regresión lineal modelo de regresión congelado
% \RegModeloOtro                  = Perceptrón multicapa otro candidato evaluado en validación
% \RegMomento                     = programacion momento de la predicción de la regresión
% \RegObjetivo                    = muestras\_dia columna objetivo de la regresión
% \RegOptimo                      = modelo = lineal, escalado = minmax, grupo = ampliado, objetivo\_log = no, seleccion = 8, codificacion = one\_hot óptimo aditivo del L16
% \RegOrigenConfiguracion         = corrida 8 del L16 de dónde sale la configuración que pasa al afinamiento
% \RegParRedundante               = media\_movil\_5 y media\_movil\_20 par de variables numéricas más correlacionado
% \RegPendienteSimple             = 1.114        pendiente de la recta simple
% \RegPorEmpate                   = no           el final se eligió por la regla de empate
% \RegPuestoConfirmacion          = 12           puesto de la confirmación entre las combinaciones del factorial
% \RegRDosSimple                  = 0.355        R cuadrado de la recta simple en entrenamiento
% \RegRDosValidacion              = 0.367        R cuadrado del modelo final en validación
% \RegRMSEConfirmacion            = 13.61        RMSE medio de la corrida de confirmación
% \RegRMSEDummy                   = 19.40        RMSE medio por bloque del Dummy (media)
% \RegRMSELinealDidactico         = 13.31        RMSE medio de la regresión lineal múltiple didáctica
% \RegRMSEMLPDidactico            = 14.31        RMSE medio del perceptrón didáctico
% \RegRMSEMejorCorrida            = 13.38        RMSE medio de la mejor corrida
% \RegRMSEMejorDidactico          = 13.31        su RMSE medio por bloque
% \RegRMSEPodado                  = 16.20        RMSE medio del árbol podado
% \RegRMSEReferencia              = 18.76        RMSE medio por bloque de la mejor referencia
% \RegRMSEReferenciaMax           = 22.51        RMSE máximo por bloque de la mejor referencia
% \RegRMSEReferenciaMin           = 15.67        RMSE mínimo por bloque de la mejor referencia
% \RegRMSESinPoda                 = 17.33        RMSE medio del árbol sin poda
% \RegRMSEValidacion              = 14.59        RMSE del modelo final en validación
% \RegRMSEValidacionInf           = 12.23        límite inferior del IC 95 % del RMSE en validación
% \RegRMSEValidacionOtro          = 15.51        RMSE del otro candidato en validación
% \RegRMSEValidacionSup           = 17.21        límite superior del IC 95 % del RMSE en validación
% \RegRejillaLinealRMSE           = 13.38        RMSE medio por bloque de la mejor combinación de la rejilla de lineal
% \RegRejillaMLPRMSE              = 14.51        RMSE medio por bloque de la mejor combinación de la rejilla de mlp
% \RegResiduoMedioValidacion      = 1.85         residuo medio (observado - predicho) en validación
% \RegRhoCompleto                 = 0.94         correlación de Spearman entre deltas del L16 y del factorial completo
% \RegSNMejorCompleto             = $-$22.57     mejor S/N (dB) del factorial completo
% \RegSNMejorCorrida              = $-$22.63     S/N (dB) de la mejor corrida
% \RegSNObservada                 = $-$22.78     S/N (dB) observada en la corrida de confirmación
% \RegSNPredicha                  = $-$22.76     S/N (dB) predicha para el óptimo aditivo
% \RegSegundosCompleto            = 7.8          segundos del factorial completo
% \RegSegundosLSixteen            = 1.0          segundos de las 16 corridas del L16
% \RegToleranciaEmpatePct         = 1            tolerancia de empate (%) de la regla de elección
% \RegToleranciaOperativa         = 10.8         tolerancia práctica en muestras por día
% \RegToleranciaOperativaPct      = 20           tolerancia práctica (% de la media diaria)
% \RegToleranciaSN                = 0.74         tolerancia (dB) de la confirmación
% \RegVariableSimple              = puntos\_programados\_dia variable de la regresión lineal simple
% \RegVariablesModelo             = 25           variables que ve el modelo congelado (tras codificar)
% --- tablas de 03_cribado_regresion ---
% tabla_catalogo.tex (16 filas): Catálogo de análisis candidatos
% tabla_cribado_agrupamiento.tex (3 filas): Cribado de agrupamiento (entrenamiento)
% tabla_cribado_clasificacion.tex (5 filas): Cribado de clasificación (entrenamiento)
% tabla_cribado_regresion.tex (5 filas): Cribado de regresión (entrenamiento)
% tabla_reg_coeficientes.tex (3 filas): Tres coeficientes estandarizados de mayor tamaño (regresión múltiple)
% tabla_reg_didactica.tex (8 filas): Pasada didáctica de regresión (media y desviación entre bloques)
% tabla_reg_l16_diseno.tex (16 filas): Diseño L16 de regresión con los niveles reales
% tabla_reg_l16_resultados.tex (16 filas): Resultados del L16 de regresión
% tabla_reg_poda.tex (6 filas): Ruta de poda del árbol de regresión
% tabla_reg_referencias.tex (3 filas): RMSE de las referencias por bloque
% tabla_reg_respuesta.tex (6 filas): Tabla de respuesta de Taguchi (regresión, S/N en dB)
% tabla_reg_validacion.tex (5 filas): Bloque de validación de la regresión con intervalos bootstrap
% tabla_tablas_modelado.tex (13 filas): Tablas de modelado construidas y auditadas
%
% ===== 04_clasificacion_bayes: 170 macros =====
% \ClfAUCCalibrado                  = 0.721        ROC-AUC de validación calibrado
% \ClfAUCConfirmacion               = 0.699        ROC-AUC medio de la corrida de confirmación
% \ClfAUCDesvMejorDidactico         = 0.021        desviación de su ROC-AUC entre bloques
% \ClfAUCEquivalentePredicha        = 0.725        ROC-AUC equivalente a la S/N predicha
% \ClfAUCLogisticaDidactico         = 0.702        ROC-AUC medio de la regresión logística didáctica
% \ClfAUCMLPDidactico               = 0.692        ROC-AUC medio del perceptrón didáctico
% \ClfAUCMejorCorrida               = 0.695        ROC-AUC medio de la mejor corrida
% \ClfAUCMejorDidactico             = 0.702        su ROC-AUC medio por bloque
% \ClfAUCPeorDidactico              = 0.634        ROC-AUC medio del más débil
% \ClfAUCPodado                     = 0.639        ROC-AUC medio del árbol podado
% \ClfAUCRejillaFinal               = 0.699        ROC-AUC medio por bloque del final en la rejilla
% \ClfAUCSinPoda                    = 0.610        ROC-AUC medio del árbol sin poda
% \ClfAUCValidacion                 = 0.721        ROC-AUC de validación del modelo que se congela
% \ClfAUCValidacionInf              = 0.697        límite inferior del IC 95 % del ROC-AUC en validación
% \ClfAUCValidacionOtro             = 0.697        ROC-AUC de validación del otro candidato
% \ClfAUCValidacionSinCalibrar      = 0.724        ROC-AUC de validación del final (antes de calibrar)
% \ClfAUCValidacionSup              = 0.745        límite superior del IC 95 % del ROC-AUC en validación
% \ClfAhorroPct                     = 19           ahorro (%) frente a la mejor política sin modelo
% \ClfAlertasSemana                 = 31.2         alertas por semana con el umbral elegido
% \ClfAlertasValidacion             = 843          alertas totales en validación
% \ClfAlfaPoda                      = 0.000122     ccp_alpha elegido para el árbol didáctico
% \ClfAnalisis                      = no\_conforme\_recepcion análisis de clasificación modelado a fondo
% \ClfAudAUCVPPEnsayo               = 0.559        ROC-AUC del VPP por ensayo frente a la contaminación real
% \ClfAudAUCVPPx                    = 0.779        ROC-AUC del VPP(x) frente a la contaminación real (validación)
% \ClfAudBrierVPPEnsayo             = 0.1742       Brier del VPP por ensayo frente a la contaminación real
% \ClfAudBrierVPPx                  = 0.1449       Brier del VPP(x) frente a la contaminación real
% \ClfAudEnsayosDentroIC            = 5            ensayos cuyo VPP real cae en el intervalo estimado
% \ClfAudFalsoRealReensayados       = 0.241        fracción realmente no contaminada entre los reensayados
% \ClfAudSesgoEspec                 = 0.0033       exceso medio de la especificidad de controles sobre la efectiva
% \ClfAudVPPReal                    = 0.756        VPP real del desarrollo (verdad de la simulación)
% \ClfBrierCalibrado                = 0.0752       Brier de validación calibrado
% \ClfBrierMejorDidactico           = 0.0754       su Brier medio
% \ClfBrierSinCalibrar              = 0.2269       Brier de validación sin calibrar
% \ClfCompletoEjecutado             = no           se ejecutó el factorial completo
% \ClfConfiguracionFinal            = modelo = logistica, escalado = estandar, grupo = ampliado, desbalance = sobremuestreo, seleccion = 26, codificacion = ordinal configuración congelada
% \ClfConfirmaAditivo               = sí           la confirmación cae dentro de la tolerancia
% \ClfControlesDesarrollo           = 20552        controles de calidad del desarrollo
% \ClfCostoAlertarTodo              = 5964         costo en validación de alertar todas
% \ClfCostoNoAlertar                = 2925         costo en validación de no alertar nunca
% \ClfDeltaPrincipal                = 1.030        delta (dB) del factor más influyente
% \ClfDeltaSegundo                  = 0.292        delta (dB) del segundo factor
% \ClfDummyAUC                      = 0.500        ROC-AUC medio del Dummy por bloque
% \ClfDummyBrier                    = 0.0813       Brier medio del Dummy por bloque
% \ClfDummyPrecisionMedia           = 0.089        precisión media del Dummy por bloque
% \ClfEnsayoEjemplo                 = mesófilos aerobios ensayo de los dos ejemplos de VPP(x)
% \ClfEnsayos                       = 7            ensayos con sensibilidad y especificidad estimadas
% \ClfEspecMax                      = 0.995        especificidad máxima entre ensayos
% \ClfEspecMin                      = 0.976        especificidad mínima entre ensayos
% \ClfEspecMinEnsayo                = mohos y levaduras ensayo con la menor especificidad
% \ClfFN                            = 353          falsos negativos en validación
% \ClfFP                            = 611          falsos positivos en validación
% \ClfFUnoUmbral                    = 0.325        F1 con el umbral elegido
% \ClfFactorPrincipal               = modelo       factor más influyente (mayor delta)
% \ClfFactorSegundo                 = grupo        segundo factor más influyente
% \ClfFalsoRiesgoAlto               = 0.026        P(falso positivo) en la muestra de riesgo alto
% \ClfFalsoRiesgoBajo               = 0.378        P(falso positivo) en la muestra de riesgo bajo
% \ClfFalsosPorCien                 = 22           falsos positivos por cada cien no conformes (controles)
% \ClfFilasReajuste                 = 29012        filas del reajuste final (entrenamiento + validación)
% \ClfHiperparametros               = C = 10       hiperparámetros congelados
% \ClfHojaRiesgoFilas               = 22           muestras en la hoja de mayor riesgo del árbol
% \ClfHojaRiesgoFraccion            = 0.773        fracción de no conformes en esa hoja
% \ClfHojaRiesgoVeces               = 9.0          veces la tasa general
% \ClfHojasPodado                   = 52           hojas del árbol podado
% \ClfHojasSinPoda                  = 546          hojas del árbol de clasificación sin poda
% \ClfInteraccionesLibres           = 9            interacciones que quedan libres en el L16
% \ClfKSeleccion                    = 26           variables que conserva la selección por F
% \ClfLogisticaParametros           = 136          coeficientes e intercepto de la regresión logística
% \ClfMLPActivacion                 = relu         función de activación de las capas ocultas
% \ClfMLPAlfa                       = 0.0010       penalización L2 (alpha) del perceptrón
% \ClfMLPCapasOcultas               = 1            número de capas ocultas del perceptrón
% \ClfMLPEntradas                   = 135          entradas del perceptrón (columnas tras el preprocesamiento)
% \ClfMLPEpocas                     = 32           épocas hasta la parada temprana
% \ClfMLPExactitudMejor             = 0.921        mejor exactitud en la fracción apartada
% \ClfMLPFraccionValidacionPct      = 15           porcentaje del entrenamiento apartado
% \ClfMLPMejorEpoca                 = 23           época con la mejor exactitud en la fracción apartada
% \ClfMLPOcultas                    = 32           neuronas por capa oculta (hidden_layer_sizes)
% \ClfMLPOptimizador                = adam         optimizador del perceptrón en scikit-learn
% \ClfMLPPaciencia                  = 8            épocas sin mejora antes de detenerse (n_iter_no_change)
% \ClfMLPParametros                 = 4385         pesos y sesgos del perceptrón
% \ClfMLPPerdidaFinal               = 0.222        pérdida logarítmica de entrenamiento en la última época
% \ClfMLPTopeEpocas                 = 150          tope de épocas (max_iter)
% \ClfMLPVecesParametros            = 32           cuántas veces más parámetros tiene la red
% \ClfMejorCorrida                  = 8            número de la mejor corrida observada del L16
% \ClfMejorDidactico                = Regresión logística mejor modelo de la pasada didáctica
% \ClfModeloFinal                   = Regresión logística modelo de clasificación congelado
% \ClfModeloMejorNivel              = regresión logística mejor nivel del factor modelo en el L16
% \ClfModeloOtro                    = Árbol de decisión otro candidato evaluado en validación
% \ClfModeloPeorNivel               = perceptrón multicapa peor nivel del factor modelo en el L16
% \ClfMomento                       = recepcion    momento de la predicción de la clasificación
% \ClfMuestrasEntrenamiento         = 22463        muestras de entrenamiento
% \ClfMuestrasValidacion            = 6549         muestras de validación
% \ClfObjetivo                      = no\_conforme columna objetivo de la clasificación
% \ClfOddsBaja                      = 0.36         su razón de odds
% \ClfOddsSube                      = 3.65         su razón de odds
% \ClfOptimo                        = modelo = logistica, escalado = estandar, grupo = ampliado, desbalance = sobremuestreo, seleccion = 26, codificacion = ordinal óptimo aditivo del L16
% \ClfOrigenConfiguracion           = corrida de confirmación de dónde sale la configuración que pasa al afinamiento
% \ClfPeorDidactico                 = KNN          modelo más débil de la pasada didáctica
% \ClfPorEmpate                     = no           el final se eligió por la regla de empate
% \ClfPositivosEntrenamiento        = 1920         no conformes en entrenamiento
% \ClfPositivosValidacion           = 585          no conformes en validación
% \ClfPositivosValidacionBayes      = 585          positivos de validación con VPP(x)
% \ClfPosteriorMasMasMax            = 0.999        posterior máxima tras dos positivos
% \ClfPosteriorMasMasMin            = 0.988        posterior mínima tras dos positivos
% \ClfPosteriorMasMenosMax          = 0.474        posterior máxima tras positivo y negativo
% \ClfPosteriorMasMenosMin          = 0.059        posterior mínima tras positivo y negativo
% \ClfPrecisionMediaMejorDidactico  = 0.254        su precisión media
% \ClfPrecisionMediaValidacion      = 0.245        precisión media en validación
% \ClfPrecisionUmbral               = 0.275        precisión con el umbral elegido
% \ClfPrevalenciaEntrenamientoPct   = 8.5          no conformes (%) en entrenamiento
% \ClfPrevalenciaValidacionPct      = 8.9          no conformes (%) en validación
% \ClfProbMediaCalibrado            = 0.099        probabilidad media en validación calibrada
% \ClfProbMediaSinCalibrar          = 0.464        probabilidad media en validación sin calibrar
% \ClfRecallUmbral                  = 0.397        recall con el umbral elegido
% \ClfRecorteCeroPct                = 0.0          positivos con previa recortada a cero (%)
% \ClfReensayados                   = 2186         positivos reensayados en el desarrollo
% \ClfReensayoConformeEsperadoPct   = 25.0         reensayos conformes esperados (%) según los controles
% \ClfReensayoConformeObservadoPct  = 27.2         reensayos conformes observados (%) en el desarrollo
% \ClfReglaBorde                    = 0.22         borde de log10(recuento / límite) de la regla de reensayo
% \ClfReglaCapturaPct               = 51           reensayos conformes que concentra el grupo (%)
% \ClfReglaGrupoPct                 = 17           reensayados con recuento que caen bajo el borde (%)
% \ClfReglaTasaAltoPct              = 15           reensayos conformes (%) sobre el borde (desarrollo)
% \ClfReglaTasaBajoPct              = 78           reensayos conformes (%) bajo el borde (desarrollo)
% \ClfReglaVecesLimite              = 1.66         veces el límite bajo las que se reensaya
% \ClfRejillaArbolAUC               = 0.654        ROC-AUC medio por bloque de la mejor combinación de la rejilla de arbol
% \ClfRejillaLogisticaAUC           = 0.699        ROC-AUC medio por bloque de la mejor combinación de la rejilla de logistica
% \ClfRevisadasPorDetectado         = 2.6          conformes revisadas por no conforme detectado
% \ClfSNMejorCorrida                = $-$3.17      S/N (dB) de la mejor corrida
% \ClfSNModeloMejor                 = $-$3.272     S/N media (dB) del mejor modelo en el L16
% \ClfSNModeloPeor                  = $-$4.302     S/N media (dB) del peor modelo en el L16
% \ClfSNObservada                   = $-$3.117     S/N (dB) observada en la corrida de confirmación
% \ClfSNPredicha                    = $-$2.798     S/N (dB) predicha para el óptimo aditivo
% \ClfSegundosCompletoEstimados     = 208          segundos estimados del factorial completo
% \ClfSegundosLSixteen              = 26.0         segundos de las 16 corridas del L16
% \ClfSemanasValidacion             = 27           semanas del periodo de validación
% \ClfSensMax                       = 0.966        sensibilidad máxima entre ensayos
% \ClfSensMin                       = 0.897        sensibilidad mínima entre ensayos
% \ClfSensMinEnsayo                 = Salmonella spp. ensayo con la menor sensibilidad
% \ClfSinRazonConformes             = 0.338        fracción conforme en ensayos sin recuento
% \ClfSinRazonReensayados           = 346          reensayados de ensayos sin recuento
% \ClfSobreDeltaAUC                 = $-$0.080     cambio de su ROC-AUC con sobremuestreo
% \ClfSobreDeltaAUCMenor            = $-$0.003     cambio de su ROC-AUC con sobremuestreo
% \ClfSobreMasPierde                = Perceptrón multicapa modelo al que el sobremuestreo le quita más ROC-AUC
% \ClfSobreMenosPierde              = Regresión logística modelo que casi no cambia con sobremuestreo
% \ClfSobremuestrea                 = sí           la tubería elegida sobremuestrea
% \ClfSorteos                       = 4000         sorteos de las posteriores Beta para el VPP
% \ClfTN                            = 5353         verdaderos negativos en validación
% \ClfTP                            = 232          verdaderos positivos en validación
% \ClfTasaMayoritaria               = 0.915        exactitud de responder siempre la clase mayoritaria
% \ClfTasaTramos                    = 0.089        tasa real de no conformes en los tramos de validación de los bloques
% \ClfToleranciaSN                  = 0.683        tolerancia (dB) de la confirmación
% \ClfUmbral                        = 0.167        umbral de alerta elegido por costo en validación
% \ClfUmbralCosto                   = 2376         costo en validación con el umbral elegido
% \ClfUmbralCostoInf                = 2196         límite inferior del IC 95 % del costo
% \ClfUmbralCostoSup                = 2542         límite superior del IC 95 % del costo
% \ClfUmbralTeorico                 = 0.167        umbral teórico por costo, CFP / (CFP + CFN)
% \ClfUmbralTeoricoCosto            = 2376         costo en validación con el umbral teórico
% \ClfUsaCalibrado                  = sí           se usa la versión calibrada
% \ClfVPPGlobal                     = 0.784        VPP del laboratorio ponderado por positivos (desarrollo)
% \ClfVPPMejor                      = 0.897        VPP del ensayo más fiable
% \ClfVPPMejorEnsayo                = Salmonella spp. ensayo con el mayor VPP
% \ClfVPPPeor                       = 0.627        VPP del ensayo menos fiable
% \ClfVPPPeorEnsayo                 = esterilidad  ensayo con el menor VPP
% \ClfVPPPeorInf                    = 0.428        límite inferior del intervalo del VPP menos fiable
% \ClfVPPPeorSup                    = 0.759        límite superior del intervalo del VPP menos fiable
% \ClfVPPxMax                       = 0.995        VPP(x) máximo entre los positivos de validación
% \ClfVPPxMin                       = 0.040        VPP(x) mínimo entre los positivos de validación
% \ClfVPPxPDiez                     = 0.597        percentil 10 del VPP(x)
% \ClfVPPxPNoventa                  = 0.968        percentil 90 del VPP(x)
% \ClfVariableBaja                  = ensayo\_Salmonella spp. variable de mayor peso que reduce el riesgo
% \ClfVariableSube                  = tipo\_producto\_lácteo variable de mayor peso que aumenta el riesgo
% --- tablas de 04_clasificacion_bayes ---
% tabla_bayes_auditoria.tex (7 filas): Auditoría del VPP frente a la verdad de la simulación
% tabla_bayes_reensayo.tex (7 filas): Posterior tras el reensayo y reensayos conformes (desarrollo)
% tabla_bayes_sens_espec.tex (7 filas): Sensibilidad y especificidad por ensayo (intervalos de Wilson)
% tabla_bayes_tramos_reensayo.tex (6 filas): Reensayos conformes por sextil de log10(recuento / límite)
% tabla_bayes_vpp_ensayo.tex (7 filas): VPP por ensayo en el desarrollo
% tabla_bayes_vpp_muestra.tex (7 filas): VPP(x) entre los positivos de validación, por ensayo
% tabla_clf_didactica.tex (5 filas): Pasada didáctica de clasificación (media y desviación entre bloques)
% tabla_clf_l16_diseno.tex (16 filas): Diseño L16 de clasificación con los niveles reales
% tabla_clf_l16_resultados.tex (16 filas): Resultados del L16 de clasificación
% tabla_clf_odds.tex (8 filas): Ocho variables con mayor peso en la regresión logística
% tabla_clf_poda.tex (6 filas): Ruta de poda del árbol de clasificación
% tabla_clf_reglas_arbol.tex (6 filas): Hojas de mayor riesgo del árbol podado
% tabla_clf_respuesta.tex (6 filas): Tabla de respuesta de Taguchi (clasificación, S/N en dB)
% tabla_clf_sobremuestreo.tex (4 filas): Modelos didácticos sin y con sobremuestreo
% tabla_clf_umbrales.tex (26 filas): Recorrido de umbrales en validación
% tabla_clf_validacion.tex (3 filas): Candidatos de clasificación en validación (intervalos bootstrap)
%
% ===== 05_agrupamiento: 74 macros =====
% \AgrAudARIDbscan                  = 0.007        ARI entre el segmento latente y DBSCAN
% \AgrAudARIKmeans                  = 0.190        ARI entre el segmento latente y K-means
% \AgrAudClientes                   = 186          clientes auditados contra el segmento latente
% \AgrAudPureza                     = 0.538        pureza de los segmentos de K-means
% \AgrAudSegmentosLatentes          = 4            segmentos latentes de la simulación
% \AgrCompARIKmDb                   = 0.051        ARI entre K-means y DBSCAN con ruido (complementario)
% \AgrCompARIKmDbSinRuido           = 0.000        ARI entre K-means y DBSCAN sin ruido (complementario)
% \AgrCompARISinEscalar             = 0.000        ARI entre K-means sin escalar y escalado (complementario)
% \AgrCompConLog                    = 1            variables con log1p en el análisis complementario
% \AgrCompDBCombinaciones           = 16           combinaciones de la rejilla de DBSCAN (complementario)
% \AgrCompDBConGrupos               = 0            combinaciones de DBSCAN con dos o más grupos (complementario)
% \AgrCompDBCriterio                = respaldo     criterio con que se eligió DBSCAN (complementario)
% \AgrCompDBEps                     = 3.42         eps elegido de DBSCAN (complementario)
% \AgrCompDBGrupos                  = 1            grupos de DBSCAN (complementario)
% \AgrCompDBMinMuestras             = 12           min_samples elegido de DBSCAN (complementario)
% \AgrCompDBRuidoPct                = 7.5          casos de ruido (%) de DBSCAN (complementario)
% \AgrCompDBSilueta                 = ---          silueta sin ruido de DBSCAN (complementario)
% \AgrCompDominante                 = antigüedad   variable que domina la varianza sin escalar (complementario)
% \AgrCompDominantePct              = 100.0        varianza (%) que aporta sin escalar (complementario)
% \AgrCompEstabilidadMax            = 0.948        ARI máximo entre semillas (complementario)
% \AgrCompEstabilidadMin            = 0.934        ARI mínimo entre semillas (complementario)
% \AgrCompFilas                     = 186          filas del análisis complementario
% \AgrCompGrupos                    = Más no conformidad y más entregas tardías, Más muestreo en planta y más tamaño de lote, Menos volumen mensual y menos muestreo en planta nombres de los segmentos de clientes
% \AgrCompHopkins                   = 0.710        Hopkins de la matriz escalada (complementario)
% \AgrCompK                         = 3            k elegido por silueta (complementario)
% \AgrCompPCAVarianzaPct            = 40.2         varianza (%) que explican las dos componentes (complementario)
% \AgrCompSegmentoMayor             = Más muestreo en planta y más tamaño de lote segmento de clientes con más muestras
% \AgrCompSegmentoMayorClientesPct  = 39.2         sus clientes (%)
% \AgrCompSegmentoMayorMuestrasPct  = 73.6         sus muestras (%)
% \AgrCompSilueta                   = 0.186        silueta de K-means con el k elegido (complementario)
% \AgrCompSiluetaMax                = 0.186        silueta máxima de la exploración (complementario)
% \AgrCompSiluetaMin                = 0.124        silueta mínima de la exploración (complementario)
% \AgrCompVariables                 = 12           variables de la matriz del análisis complementario
% \AgrComplementario                = segmentos\_clientes análisis de agrupamiento complementario
% \AgrComplementarioSupera          = sí           el complementario habría superado al seleccionado
% \AgrEfectoFilas                   = 0.001        cambio de silueta por las filas
% \AgrEfectoLog                     = $-$0.286     cambio de silueta por log1p
% \AgrKMax                          = 10           mayor k explorado en K-means
% \AgrKMin                          = 2            menor k explorado en K-means
% \AgrPuntajeCompEtapa              = 0.124        puntaje de la regla del complementario con la silueta de la etapa 05
% \AgrPuntajeSelEtapa               = 0.049        puntaje de la regla del seleccionado con la silueta de la etapa 05
% \AgrRuidoMaximoPct                = 15           ruido máximo (%) de una solución DBSCAN con grupos
% \AgrRuidoRespaldoPct              = 10           ruido máximo (%) de DBSCAN como detector de atípicos
% \AgrSelARIKmDb                    = 0.135        ARI entre K-means y DBSCAN con ruido (seleccionado)
% \AgrSelARIKmDbSinRuido            = 0.064        ARI entre K-means y DBSCAN sin ruido (seleccionado)
% \AgrSelARISinEscalar              = 0.000        ARI entre K-means sin escalar y escalado (seleccionado)
% \AgrSelConLog                     = 5            variables con log1p en el análisis seleccionado
% \AgrSelDBCombinaciones            = 16           combinaciones de la rejilla de DBSCAN (seleccionado)
% \AgrSelDBConGrupos                = 10           combinaciones de DBSCAN con dos o más grupos (seleccionado)
% \AgrSelDBCriterio                 = regla        criterio con que se eligió DBSCAN (seleccionado)
% \AgrSelDBEps                      = 1.59         eps elegido de DBSCAN (seleccionado)
% \AgrSelDBGrupos                   = 2            grupos de DBSCAN (seleccionado)
% \AgrSelDBMinMuestras              = 16           min_samples elegido de DBSCAN (seleccionado)
% \AgrSelDBRuidoPct                 = 12.7         casos de ruido (%) de DBSCAN (seleccionado)
% \AgrSelDBSilueta                  = 0.314        silueta sin ruido de DBSCAN (seleccionado)
% \AgrSelDominante                  = volumen de muestra variable que domina la varianza sin escalar (seleccionado)
% \AgrSelDominantePct               = 95.0         varianza (%) que aporta sin escalar (seleccionado)
% \AgrSelEstabilidadMax             = 0.996        ARI máximo entre semillas (seleccionado)
% \AgrSelEstabilidadMin             = 0.981        ARI mínimo entre semillas (seleccionado)
% \AgrSelFilas                      = 5000         filas del análisis seleccionado
% \AgrSelGrupos                     = Recepción tardía y transporte más frío, Transporte largo, Recepción temprana, Plazo prometido largo y entrega lenta nombres de los perfiles de muestra
% \AgrSelHopkins                    = 0.714        Hopkins de la matriz escalada (seleccionado)
% \AgrSelK                          = 4            k elegido por silueta (seleccionado)
% \AgrSelPCAVarianzaPct             = 38.1         varianza (%) que explican las dos componentes (seleccionado)
% \AgrSelSilueta                    = 0.148        silueta de K-means con el k elegido (seleccionado)
% \AgrSelSiluetaMax                 = 0.148        silueta máxima de la exploración (seleccionado)
% \AgrSelSiluetaMin                 = 0.120        silueta mínima de la exploración (seleccionado)
% \AgrSelVariables                  = 8            variables de la matriz del análisis seleccionado
% \AgrSeleccionado                  = perfiles\_muestra análisis de agrupamiento seleccionado por la regla
% \AgrSemillasEstabilidad           = 5            semillas adicionales para la estabilidad
% \AgrSiluetaCribado                = 0.433        silueta del seleccionado en el cribado
% \AgrSiluetaSinLog                 = 0.434        silueta con filas de la etapa 05 sin log1p
% \AgrToleranciaSilueta             = 0.02         diferencia de silueta que se discute
% \AgrUmbralNombre                  = 0.5          desviaciones estándar para nombrar un grupo
% --- tablas de 05_agrupamiento ---
% tabla_agr_comp_grupos.tex (3 filas): Grupos de K-means en segmentos_clientes con su nombre
% tabla_agr_comp_kmeans.tex (9 filas): Inercia y silueta por k (segmentos_clientes)
% tabla_agr_comp_medias.tex (12 filas): Medias por grupo en segmentos_clientes (variables originales)
% tabla_agr_comp_segmentos.tex (3 filas): Tamaño y desempeño de cada segmento de clientes
% tabla_agr_sel_grupos.tex (4 filas): Grupos de K-means en perfiles_muestra con su nombre
% tabla_agr_sel_kmeans.tex (9 filas): Inercia y silueta por k (perfiles_muestra)
% tabla_agr_sel_medias.tex (8 filas): Medias por grupo en perfiles_muestra (variables originales)
% tabla_agr_silueta_cribado.tex (2 filas): Silueta y puntaje del cribado con el preprocesamiento de cada etapa
%
% ===== 06_evaluacion_final: 77 macros =====
% \FinAUCPrueba                        = 0.724        ROC-AUC en la prueba
% \FinAUCPruebaInf                     = 0.705        límite inferior del IC 95 % del ROC-AUC en la prueba
% \FinAUCPruebaSup                     = 0.743        límite superior del IC 95 % del ROC-AUC en la prueba
% \FinAgrHopkins                       = 0.714        Hopkins de los perfiles de muestra
% \FinAgrK                             = 4            perfiles de muestra de K-means
% \FinAgrSegmentos                     = 3            segmentos de clientes de K-means
% \FinAgrSilueta                       = 0.148        silueta de los perfiles de muestra (etapa 05)
% \FinAgrSiluetaCribado                = 0.433        silueta del mismo análisis en el cribado
% \FinAhorroPct                        = 15           ahorro (%) frente a la mejor política sin modelo en la prueba
% \FinAlertasSemana                    = 21.2         alertas por semana en la prueba
% \FinAudBrierVPPEnsayo                = 0.1586       Brier del VPP por ensayo en la prueba
% \FinAudBrierVPPx                     = 0.1381       Brier del VPP(x) en la prueba
% \FinAudVPPReal                       = 0.800        fracción de positivos de prueba realmente contaminados
% \FinBrierPrueba                      = 0.0753       Brier en la prueba
% \FinCambioRMSEPct                    = $-$15        cambio (%) del RMSE de validación a prueba
% \FinClfModelo                        = Regresión logística modelo de clasificación evaluado en la prueba
% \FinCostoAlertarTodo                 = 7882         costo en la prueba de alertar todo
% \FinCostoNoAlertar                   = 3875         costo en la prueba de no alertar
% \FinCostoPorMuestraPrueba            = 0.381        costo por muestra en la prueba
% \FinCostoPrueba                      = 3295         costo con el umbral en la prueba
% \FinCostoRetro                       = 3231         costo con el umbral retrospectivo (no se usa)
% \FinDiasFinDeAnio                    = 5            días hábiles de prueba en la ventana de fin de año
% \FinDiasFueraPct                     = 37.9         días de prueba fuera de tolerancia (%), modelo
% \FinDiasFueraReferenciaPct           = 57.1         días de prueba fuera de tolerancia (%), referencia
% \FinDiasPrueba                       = 203          días hábiles de prueba
% \FinEnsayosPocosPositivos            = 2            ensayos con menos de 30 positivos de prueba
% \FinFN                               = 523          falsos negativos en la prueba
% \FinFP                               = 680          falsos positivos en la prueba
% \FinFUnoPrueba                       = 0.295        F1 con el umbral en la prueba
% \FinHabilidadPrueba                  = 0.334        habilidad del modelo en la prueba
% \FinHabilidadPruebaInf               = 0.247        límite inferior del IC 95 % de la habilidad en la prueba
% \FinHabilidadPruebaSup               = 0.408        límite superior del IC 95 % de la habilidad en la prueba
% \FinMAEPrueba                        = 9.80         MAE del modelo de regresión en la prueba
% \FinMediaPrueba                      = 44.1         muestras por día hábil en promedio (prueba)
% \FinMinimoPositivos                  = 30           mínimo de positivos de prueba por ensayo para una cifra precisa
% \FinMuestrasPrueba                   = 8657         muestras de prueba con resultado
% \FinPositivosPrueba                  = 775          positivos (no conformes) en la prueba
% \FinPositivosReensayados             = 691          positivos de la prueba de clasificación (recibidos en la prueba) que tienen reensayo
% \FinPrecisionMediaPrueba             = 0.243        precisión media en la prueba
% \FinPrecisionPrueba                  = 0.270        precisión con el umbral en la prueba
% \FinPrecisionPruebaInf               = 0.242        límite inferior del IC 95 % de la precisión
% \FinPrecisionPruebaSup               = 0.299        límite superior del IC 95 % de la precisión
% \FinPrevalenciaPrueba                = 0.090        tasa real de no conformes en la prueba
% \FinProbMediaPrueba                  = 0.090        probabilidad media predicha en la prueba
% \FinPruebaDesde                      = 2026-01-02   primer día de prueba
% \FinPruebaHasta                      = 2026-10-30   último día de prueba
% \FinRDosPrueba                       = 0.529        R cuadrado del modelo de regresión en la prueba
% \FinRMSEPrueba                       = 12.47        RMSE del modelo de regresión en la prueba
% \FinRMSEPruebaInf                    = 11.20        límite inferior del IC 95 % del RMSE en la prueba
% \FinRMSEPruebaSup                    = 13.66        límite superior del IC 95 % del RMSE en la prueba
% \FinRecallPrueba                     = 0.325        recall con el umbral en la prueba
% \FinRecallPruebaInf                  = 0.293        límite inferior del IC 95 % del recall
% \FinRecallPruebaSup                  = 0.357        límite superior del IC 95 % del recall
% \FinReensayoConformeEsperadoPct      = 24.1         reensayos conformes esperados (%) en la prueba
% \FinReensayoConformeObservadoInfPct  = 20.3         límite inferior de Wilson (%)
% \FinReensayoConformeObservadoPct     = 23.3         reensayos conformes observados (%) en la prueba
% \FinReensayoConformeObservadoSupPct  = 26.6         límite superior de Wilson (%)
% \FinReensayosLecturaPrueba           = 698          reensayos cuya lectura cae en la prueba (la tabla de reensayo se parte por fecha de lectura)
% \FinReensayosRecibidosAntes          = 7            reensayos leídos en la prueba de muestras recibidas antes del corte (fuera de la prueba de clasificación)
% \FinRegModelo                        = Regresión lineal modelo de regresión evaluado en la prueba
% \FinRegReferencia                    = base\_media\_5 referencia de la habilidad en la prueba
% \FinReglaBajoBorde                   = 93           reensayos de prueba bajo el borde
% \FinReglaCapturaPct                  = 55           reensayos conformes que concentra el grupo en la prueba (%)
% \FinReglaReensayosConRecuento        = 593          reensayos de prueba con recuento
% \FinReglaTasaAltoPct                 = 12           reensayos conformes (%) sobre el borde en la prueba
% \FinReglaTasaBajoInfPct              = 68           límite inferior de Wilson (%)
% \FinReglaTasaBajoPct                 = 77           reensayos conformes (%) bajo el borde en la prueba
% \FinReglaTasaBajoSupPct              = 85           límite superior de Wilson (%)
% \FinResiduoFinDeAnio                 = $-$12.5      residuo medio en la ventana de fin de año
% \FinResiduoResto                     = 0.6          residuo medio fuera de la ventana de fin de año
% \FinSemanasPrueba                    = 44           semanas del periodo de prueba
% \FinTN                               = 7202         verdaderos negativos en la prueba
% \FinTP                               = 252          verdaderos positivos en la prueba
% \FinUmbral                           = 0.167        umbral congelado de la clasificación
% \FinUmbralRetro                      = 0.140        umbral retrospectivo de menor costo en la prueba (no se usa)
% \FinUnoDeCadaFalso                   = 5            uno de cada cuántos positivos sería falso (prueba)
% \FinVPPxMedio                        = 0.794        VPP(x) medio de los positivos de prueba
% --- tablas de 06_evaluacion_final ---
% tabla_final_auditoria.tex (7 filas): Auditoría en prueba del VPP frente a la verdad
% tabla_final_bayes.tex (7 filas): VPP y reensayos conformes en la prueba, por ensayo
% tabla_final_clasificacion_validacion.tex (7 filas): Validación frente a prueba, clasificación
% tabla_final_congelado.tex (2 filas): Modelos congelados antes de abrir la prueba
% tabla_final_preguntas.tex (4 filas): Pregunta del laboratorio, respuesta, evidencia y limitación
% tabla_final_regla_reensayo.tex (2 filas): Regla de reensayo en desarrollo y prueba
% tabla_final_regresion.tex (3 filas): Regresión en la prueba (modelo y referencias)
% tabla_final_regresion_validacion.tex (4 filas): Validación frente a prueba, regresión
% tabla_final_resumen.tex (8 filas): Resumen final: validación (o desarrollo) frente a prueba
%
